\input{preamble}

\begin{document}

\title{Exact Alternatives and the Price of\\Surrogate Certification
in Support-Aware Compression}

\author{\IEEEauthorblockN{Vu Phan An Nguyen}
\IEEEauthorblockA{\textit{Rider University} \\
Lawrenceville, NJ, USA}}

\maketitle

\begin{abstract}
A compression is chosen by comparing losses, and a learner choosing one under noise faces an
information game whose alternative set is the collection of worlds in which the choice would
be wrong. That set is routinely replaced by a surrogate: the extrema realising each loss are
frozen at their current \emph{indices} in the loss formula and the ordering constraints are
dropped --- which is not the same as forbidding the other cells to move, and an earlier
draft described it that way. We show that optimising the
surrogate can cost an unbounded factor, and --- the part that makes it a statement about
\emph{design} rather than about two numbers --- that the surrogate's own optimum collects no
information about the true decision at all.

Our main theorem is a separation \emph{after the allocation has been optimised}, with both
games taking their infimum over the same parameter class, at the same world, under the same
observation law: on a construction with four extremum cells and $r$ pairs of middle cells at
variance $M$, the frozen game has the exact value $9/512$ while the true game lies between
$9/128(4+2rM)$ and $1/8Mr$, so the ratio is $\Theta((rM)^{-1})$. The gap is already there at
$M=1$: it comes from adding actions, not from inflating a variance. Restricting an infimum
raises it, so neither bound is inherited from a larger set --- both witnesses are shown to
lie in the class --- and the alternative sets, though closed, are approached from inside by
strictly alternative worlds.

The consequence is the second theorem. The frozen value is at most $9/512$ times the mass on
the extremum cells, so any design attaining the surrogate's optimum puts \emph{exactly zero}
weight on the middle cells and has zero information against the alternatives that matter;
and a design merely within $\varepsilon$ of that optimum leaks at most $64\varepsilon/9rM$ ---
$\varepsilon$ times the whole separation factor. Solving the surrogate more accurately drives
the allocation towards a design that cannot distinguish the alternatives the compression
turns on.

Supporting these: the alternative set is pinned exactly, as an \emph{equality} with a finite
union of polyhedron-intersect-halfspace pieces proved in both directions, with the strict
inequality handled by an explicit modulus rather than silently closed; certificates at two
levels, a regional KKT certificate and a dual over the whole query set, so ``optimal'' is
never shorthand for ``optimal within the family we searched''; and a learner that stops under
a certificate which is a weak-duality \emph{lower} bound, never a primal value.

The separation reaches sample complexity for the design itself. The learner that tracks an
exact surrogate optimum, with forced exploration at rate $1/\sqrt{s+1}$ and the true
likelihood-ratio stop, pays $\Omega(\log^{2}(1/\delta))$ samples, while the ideal learner,
defined and proved here, pays $\Theta(\log(1/\delta))$. This holds for the deployable learner
that re-solves the surrogate at its own estimate every round, not only for the surrogate's
optimum at the true world. Its oracle sees only the estimate, and the rounds in which the
estimate misleads it cost a constant in expectation. The measured panel is reported per confidence level under common random
numbers, and supports the implementation, not the asymptotics.
\end{abstract}

\begin{IEEEkeywords}
compression, alternative sets, experimental design, convex certificates, formal verification
\end{IEEEkeywords}

\section{Introduction}
\label{sec:intro}

\subsection{The question}

Let a world be $\theta$, let $T$ index a retained coordinate set, and let $L_{\theta,K}(T)$
be the compression loss of $T$ measured on the deployment support $K$. The learner must
return a $T$ whose regret is at most $\varepsilon$, and it pays for observations. Its cost
is governed by an information game whose inner infimum ranges over the \emph{alternatives}:
the worlds in which the returned $T$ would be wrong.

Computing that infimum requires knowing the alternative set. The natural shortcut is to
freeze the argmax and argmin realising each loss and ask how far the observations must move
to reverse the comparison. This paper is about how badly that shortcut misprices the game,
and about what it takes to replace it with something certified.

\subsection{What is new, against the nearest sources}

Information games with several correct answers, and alternative sets that decompose into
finitely many convex regions, are developed in the literature; Degenne and
Koolen~\cite{degenne2019} for the multiple-answer setting, and Wang and
You~\cite{wangyou2026} for the piecewise-convex decomposition. We take that structure as
given and do not claim it. We also do not overstate it: \cite{wangyou2026} leaves an exact
stopping threshold problem-dependent and labels part of its scaling argument heuristic, so
it should not be cited as having closed the algorithmic question.

Bennouna et al.~\cite{bennouna2026} study data informativeness in linear optimisation. That
is the closest work on ``which observations matter for a decision'', and its scope is linear
optimisation; we do not extend its statements to noisy compression, which it does not make.

Two things this paper does \emph{not} claim, said first because both are easy to read into
it. It does not claim that an adaptive allocation attains the asymptotic complexity for
several correct answers: \cite{degenne2019} gives both the lower bound and an algorithm
attaining it, and we inherit that machinery. And it does not claim to open the question of
when a proxy objective suffices for a decision; sufficiency criteria for decision proxies
are studied directly, e.g.~\cite{schutte2026}. Our statement is about what the complexity
\emph{is} once the decision is compression, and how far a frozen-linear surrogate is from
it after every allocation has been optimised.

The contribution is a chain, and each link is load-bearing:
real compression $\to$ its real alternative geometry $\to$ the \emph{optimal} surrogate
design is still wrong by an unbounded factor $\to$ the cost of learning changes. The last
link is proved for learners that are defined completely, and for no others. An
\emph{ideal} learner, whose oracle value and stopping statistic are exact infima, is
$\delta$-correct and reaches the corrected benchmark. The learner that re-solves the
surrogate at its own estimate, with the same forcing and the same stop, pays a whole power of
$\log(1/\delta)$ more (Section~\ref{sec:adaptive}). The implementation of
Section~\ref{sec:learner} is a different algorithm, with finite precision and a bracketing oracle. For it the rate is proved only
conditionally, on named information-accumulation and expectation hypotheses, together with
the repair of the one hypothesis it violated outright (its threshold grows without bound in
$t$; see the shift argument below), and it is reported as a finite-panel measurement.
Concretely:
\begin{enumerate}
\item[1)] \textbf{An equality, not a decomposition found by search.} We prove the
alternative set \emph{equals} the union of its regions, and we handle the strict inequality
with an explicit modulus, including the case in which the strict part is empty and the
region must be dropped rather than closed.
\item[2)] \textbf{Certificates at both levels.} A KKT certificate for each region, and a
dual over the \emph{whole} query set for the design---so ``optimal'' is not shorthand for
``optimal within the family we searched''.
\item[3)] \textbf{Two-sided failure of the frozen surrogate}, with explicit families and
rates in both directions, so no constant-factor repair is possible; and the separation
$\Theta((rM)^{-1})$ survives at \emph{uniform} variance, so it is not bought by inflating
a variance.
\item[4)] \textbf{A consequence, not only a price gap.} Two different game values do not by
themselves make the surrogate's policy more expensive. Theorem~\ref{thm:starve} shows the
surrogate's optimum \emph{starves} the cells that carry the compression, and that an
$\varepsilon$-optimal design still does, at rate $64\varepsilon/(9Mr)$.
\item[5)] \textbf{Attainability, for an ideal learner.} A learner defined completely (a warm
start, exact-value selection over a refining net, forcing at rate $1/\sqrt{s+1}$, and the
exact generalised-likelihood-ratio stop) is $\delta$-correct, has
$\mathbb{E}\tau_{\delta}<\infty$, and satisfies
$\mathbb{E}\tau_{\delta}/\log(1/\delta)\to1/\Gamma_{\rm true}$ on the separating family, with no
hypothesis beyond $M,r,a>0$.
\item[6)] \textbf{The sample-cost separation, for the surrogate's design.} Tracking any exact
surrogate optimum, with the same forcing and the same stop, costs
$\mathbb{E}\tau_{\delta}\ge(KM/8)^{2}\log^{2}(1/\delta)$. The ratio to the ideal learner is
unbounded. This holds for the surrogate's optimum at the true world, and for the deployable
learner that re-solves the surrogate at its own estimate every round.
\end{enumerate}
Items 1)--4) are statements about the game and the design it selects. Items 5) and 6) carry
them to samples, for learners that are defined completely. None of them is the
implementation of Section~\ref{sec:learner}. The count bound that 6) needs is proved
directly: forcing and the rounds with a misleading estimate are the only sources of middle
queries, and the latter are finitely many in expectation. The chain of
Section~\ref{sec:alloc} is therefore not needed for 6). Theorems~\ref{thm:sepcentred}
and~\ref{thm:starve} are the main results on the game. The implemented learner is reported in Section~\ref{sec:evidence} as a measurement on
one finite panel, and carries no asymptotic claim.


\paragraph{Which of our theorems does not follow.}
Stated at the level of theorems rather than topics, so a reader can check the delta without
reading the history.

\begin{table}[h]
\centering
\footnotesize
\begin{tabular}{@{}p{0.30\columnwidth}p{0.62\columnwidth}@{}}
\toprule
Nearest source & Our theorem that does not follow from it \\
\midrule
Proxy insufficiency for a decision & The information game \emph{after} optimisation: the
two-sided bracket $4/(4{+}2rM)\le\Gamma_{\rm true}/\Gamma_{\rm fr}\le 64/9rM$, and that
\emph{every} surrogate optimum has $J_{\rm true}=0$ with a proportional near-optimal version \\
Multiple-answer track-and-stop & The alternative geometry of this compression loss, hence the
repair frontier $\mathcal F(\beta)=\Theta(\beta/rM)$ with matching two-sided bounds \\
Acceptance sets, cost-aware sampling & The starvation consequence for the sampling rule, and
the window head the rate must pay: $b/G+O(\sqrt b)$, not $b/G+O(1)$ \\
\bottomrule
\end{tabular}
\end{table}

\section{The model}
\label{sec:model}

Parametrise the world linearly, $F_{\theta}(x)=\phi(x)^{\top}\theta$ on a finite domain $X$
with a fixed reference and a deployment support $K\subseteq X$. With the reference fixed the
product responses are linear in $\theta$, so for a retained set $T$ the residual at a support
point is a linear form, $r_{\theta,T}(x)=a_{T}(x)^{\top}\theta$, and
\begin{equation}
L_{\theta,K}(T)=\tfrac12\Bigl[\max_{x\in K}a_{T}(x)^{\top}\theta-\min_{x\in K}a_{T}(x)^{\top}\theta\Bigr].
\label{eq:loss}
\end{equation}
Observations are Gaussian and linear in $\theta$: query $e$ returns $Y_{e}\sim N(M_{e}\theta,\sigma_{e}^{2})$
at cost $c_{e}$, and a design is a non-negative allocation $w$ with $\sum_{e}c_{e}w_{e}=1$.

\NOTEPAR{Superseded, and corrected rather than deleted. This paragraph said the two bridges
were asserted by fixing the model. They are proved. The coefficient vector $a_{T}$ is
\emph{defined} from $F$, the reference $q$ and the retained set $T$, and
\lean{residual\_eq\_inner} is the identity that makes it the residual, so it is derived and
not handed over; \lean{loss\_eq\_half\_range} is \eqref{eq:loss}; and
\lean{selectedRadius\_eq\_compressionLoss} identifies the corpus's own
\lean{selectedRadius} on the extended family with half the oscillation of that residual.
What is \emph{not} proved, and is stated where it arises rather than here, is that the
residual equals any expression with responses dropped---it does not in general.}

\section{The alternative set, exactly}
\label{sec:alt}

\subsection{The equality}

Fix $T$ and a competitor $S$. Pinning which support points realise the four extrema cuts out
a polyhedron, and on it the doubled loss difference is a single linear form.

\begin{theorem}[Region form]
\label{thm:region}
If $x,y$ realise the max and min of $r_{\theta,T}$ and $p,q$ those of $r_{\theta,S}$, then
\begin{equation}
2\bigl[L_{\theta,K}(T)-L_{\theta,K}(S)\bigr]
= \bigl(a_{T}(x)-a_{T}(y)-a_{S}(p)+a_{S}(q)\bigr)^{\top}\theta .
\label{eq:regionform}
\end{equation}
\end{theorem}

\leanbox{\lean{loss\_eq\_of\_extrema}, \lean{region\_form\_is\_linear} in
\lean{LeanPaperAAlternativesV1}}

\begin{theorem}[The alternative set is the union]
\label{thm:alteq}
$L_{\theta,K}(T)-L_{\theta,K}(S)>\varepsilon$ holds if and only if there exist extrema
indices satisfying the region conditions and making the linear form
of~\eqref{eq:regionform} exceed $2\varepsilon$.
\end{theorem}

\leanbox{\lean{alt\_iff\_exists\_region}, with \lean{exists\_extrema}, in
\lean{LeanPaperAAlternativesV1}}

Both directions matter. Left to right says the true extrema always pin \emph{some} region,
so the union is complete; right to left says every piece consists of genuine alternatives,
so nothing spurious is added. Neither is a search result.

\subsection{The closure, where care is needed}

The alternative condition is strict, and an infimum over a strict region is not in general
the infimum over its closure. The exchange is legitimate exactly when the region contains a
strictly bad point.

\begin{theorem}[Closure exchange]
\label{thm:closure}
If the region is convex and contains $\theta_{+}$ with $v^{\top}\theta_{+}>c$, then the
infimum of the energy over $\{v^{\top}\theta>c\}$ equals its infimum over
$\{v^{\top}\theta\ge c\}$.
\end{theorem}

\leanbox{\lean{strict\_inf\_eq\_closed\_inf}, with \lean{segment\_stays\_strict} and
\lean{energy\_segment\_bound}, in \lean{LeanPaperAAlternativesV1}}

The proof needs no limiting argument: along the segment from a closed point towards
$\theta_{+}$ the energy moves by at most $t$ times an explicit constant, and the segment is
strictly bad for every $t>0$. When the strict part is \emph{empty} the region contributes no
alternative and must be dropped; closing it would manufacture an alternative out of a tie.

\leanbox{\lean{strict\_empty\_has\_no\_alternatives} in the same module}

\section{Certifying, at both levels}
\label{sec:cert}

\subsection{Each region}

\begin{theorem}[KKT certificate]
\label{thm:kkt}
Let $H$ be symmetric positive semidefinite and let $\lambda$ be feasible for $A\lambda\le b$
with multipliers $y\ge0$ satisfying $H(\lambda-\theta)+A^{\top}y=0$ and complementary
slackness. Then $\lambda$ minimises $\tfrac12(\cdot-\theta)^{\top}H(\cdot-\theta)$ over the
region, against every feasible competitor.
\end{theorem}

\leanbox{\lean{kkt\_region\_minimum}, with \lean{quad\_expand}, in
\lean{LeanPaperAOracleCertificateV1}}

$H$ is \emph{not} assumed invertible. That matters: a design leaving a direction unmeasured
produces a singular $H$, and a statement requiring invertibility would silently exclude the
designs one most wants to compare.

\subsection{Every design}

\begin{theorem}[Design optimality over the whole query set]
\label{thm:dual}
Let $\nu$ be a prior on alternatives with $\sum_{\ell}\nu_{\ell}=1$, $\nu\ge0$. If
$\sum_{\ell}\nu_{\ell}d_{e}(\theta,\lambda_{\ell})\le Vc_{e}$ for \emph{every} admissible
query $e$, then $J_{T}(w,\theta)\le V$ for every design $w$ with
$\sum_{e}c_{e}w_{e}=1$. With a primal allocation attaining $V$, $\sup_{w}J_{T}(w,\theta)=V$.
\end{theorem}

\leanbox{\lean{design\_upper\_bound\_of\_dual}, \lean{design\_value\_eq\_of\_primal\_dual} in
\lean{LeanPaperAOracleCertificateV1}}

This is the step that converts ``best in the family we parametrised'' into ``best''. If a
query violates the dual inequality, that violation is a direction of improvement: add the
query and re-solve.

\section{The frozen surrogate misprices, both ways}
\label{sec:proxy}

\subsection{Two families}

Let $M$ be the effective variance of one designated cheap cell; all others have variance
one, and the decision compares $\mathrm{range}(A)$ with $\mathrm{range}(B)$.

\emph{Family 1, the surrogate prices too high.} $A=(0,1,\tfrac12)$, $B=(0,\tfrac34,\tfrac12)$
with variance $M$ on $B$'s middle cell. Freezing the extrema forces four unit-variance cells
to close a gap of $\tfrac14$, giving $J_{\rm frozen}=\tfrac1{128}$. The true alternative may
instead lift $B$'s middle cell to become the new maximum, at cost $1/8M$, so
\begin{equation}
J_{\rm full}/J_{\rm frozen}\le 16/M\longrightarrow0 .
\label{eq:fam1}
\end{equation}

\emph{Family 2, the surrogate prices too low.} $A=(0,1,\tfrac34)$, $B=(0,\tfrac12,\tfrac14)$
with variance $M$ on $A$'s maximum, so $J_{\rm frozen}=1/8(M+3)$. Lowering the cheap maximum
is not enough because the runner-up at $\tfrac34$ remains; the true cost stays bounded below
and
\begin{equation}
J_{\rm full}/J_{\rm frozen}\ \ge\ (M+3)/24\longrightarrow\infty .
\label{eq:fam2}
\end{equation}

\subsection{Computed, with every region certified}

A word on what $J_{\rm frozen}$ is, because the ratios are only coherent under the right
reading. $J_{\rm frozen}$ is the \emph{frozen-linear} surrogate: the range comparison is
linearised at the current extrema and the ordering conditions are dropped. It is \emph{not}
the minimum over the region that pins those extrema. Were it that, a region being a subset
of parameter space would force $J_{\rm full}\le J_{\rm region}$ and a ratio above one would
be impossible. Family~2's ratios exceed one because the surrogate optimises over a set that
is \emph{not} contained in the true alternative set: it admits parameters violating the
ordering conditions, which are not alternatives at all. Neither set contains the other, and
that is precisely why the surrogate can err in both directions.

\begin{table}[t]
\caption{The frozen surrogate against the true oracle. Each cell solves all $81$ regions and
every one of them carries a KKT certificate, not only the winner. \NUM{}}
\label{tab:proxy}
\centering
\footnotesize
\begin{tabular}{@{}rrrrr@{}}
\toprule
$M$ & Fam.\,1 ratio & $16/M$ & Fam.\,2 ratio & $(M{+}3)/24$ \\
\midrule
$1$ & $1.000$ & $16$ & $1.000$ & $0.167$ \\
$64$ & $0.2388$ & $0.25$ & $4.59$ & $2.79$ \\
$256$ & $0.0618$ & $0.0625$ & $16.58$ & $10.79$ \\
$1024$ & $0.0156$ & $0.0156$ & $64.58$ & $42.79$ \\
\bottomrule
\end{tabular}
\end{table}

Two features deserve emphasis. The surrogate is \emph{exact} at small $M$ and only breaks
past a threshold, so the failure is a genuine transition rather than a modelling artefact.
And because the error runs in both directions on the same decision rule, no constant-factor
guarantee relating $J_{\rm frozen}$ to $J_{\rm full}$ can exist.

\TODOPAR{These two families are stated for an oracle over cell ranges with per-cell
observations. Realising them inside the model of Section~\ref{sec:model}---an explicit
$F_{\theta}$, reference and admissible queries---has \emph{not} been done. If the move from
per-cell to per-response observation changes the Hessian, the energies must be recomputed;
the constants above must not be transported by assertion. This is the first thing the final
version owes.}

\section{Certifying a margin costs more than deciding}
\label{sec:margin}

On the full cube take $F(x_{1},x_{2})=ax_{1}+bx_{2}+cx_{1}x_{2}$ with $a>b>0$, $c>0$ and a
corner reference. Then $L(\{1\})=(b+c)/2$ and $L(\{2\})=(a+c)/2$, so the \emph{decision} is
settled by the sign of $a-b$ alone, while the interaction bound is $\delta_{I}=c$ and the
Top-$C$ margin \emph{certificate} $\rho\ge2\delta_{I}$ reads $a-b\ge2c$. Both are half-space
questions about $\theta$, but about different half-spaces.

\begin{table}[t]
\caption{Cost of the margin certificate against the decision it justifies, at $a=3$, $b=1$,
$\varepsilon=0.1$. \NUM{}}
\label{tab:margin}
\centering
\footnotesize
\begin{tabular}{@{}rrrr@{}}
\toprule
$c$ & $\Gamma_{\rm dec}$ & $\Gamma_{\rm cert}$ & ratio \\
\midrule
$0.25$ & $0.6050$ & $0.0313$ & $19.4$ \\
$0.75$ & $0.6050$ & $0.0035$ & $174.2$ \\
$0.90$ & $0.6050$ & $0.0006$ & $1089.0$ \\
$\ge1.00$ & $0.6050$ & $0$ & $\infty$ \\
\bottomrule
\end{tabular}
\end{table}

\begin{theorem}[Certifying is never cheaper]
\label{thm:cert}
If every world satisfying the certificate condition is one in which $T$ is safe, then the
decision's alternatives are contained in the certificate's, so
$\Gamma_{\rm cert}\le\Gamma_{\rm dec}$.
\end{theorem}

\leanbox{\lean{certification\_costs\_at\_least\_as\_much} in
\lean{LeanPaperAOracleCertificateV1}}

At $a-b\le2c$ the true world itself fails the margin condition, so the certificate is
unattainable at any cost while the decision remains settled.

\NOTEPAR{Two scope limits. This is an \emph{information-rate} separation, not a ratio of
expected sample counts: without a matching stopping theorem the two are not interchangeable.
And the lower bound is against \emph{this} certificate, the one the corpus uses; a weaker
sufficient condition tailored to the family might close much of the gap, and no claim is made
about ``every certificate''.}

\section{No constant-factor guarantee, after the allocation is optimised}
\label{sec:separation}

The mispricing above is at a fixed design. The question a practitioner asks is whether it
survives choosing the design well, so the statement below is about the games
\[
\Gamma_{\bullet}(\theta)\;=\;\sup_{w\in\Delta}\ \inf_{\lambda\in\mathrm{Alt}_{\bullet}
\cap\,\Theta_{0}}\ \tfrac12\sum_{e}\frac{w_{e}(\lambda_{e}-\theta_{e})^{2}}{\nu_{e}},
\]
with $\mathrm{Alt}_{\rm true}$ the worlds in which retaining $A$ is strictly worse and
$\mathrm{Alt}_{\rm proxy}$ the worlds satisfying the frozen contrast.

\paragraph{The model class, and why the tables must be centred.}
$\Theta_{0}$ is the centred class: each arm's table sums to zero. That is not decoration.
On the centred class the product response \emph{is} the table, with no additive remainder,
so the energy above is the information the queries carry rather than a stand-in for it. The
tables as printed are not centred, so a statement naming the raw cell values and asserting
$\theta\in\Theta_{0}$ has contradictory hypotheses and says nothing. An earlier version of
our wrapper did exactly that. The repair is to subtract each arm's mean and to prove that
this changes nothing the decision can see: a common shift of an arm moves neither its range,
nor the frozen contrast, nor---when the world and the candidate are shifted together---the
energy. The worlds used below are therefore
\[
\theta^{-}=\bigl(-\tfrac12,\tfrac12,0,0;\ -\tfrac7{16},\tfrac5{16},\tfrac1{16},
\tfrac1{16}\bigr),\quad
\theta^{+}=\bigl(-\tfrac9{16},\tfrac7{16},\tfrac3{16},-\tfrac1{16};\
-\tfrac14,\tfrac14,0,0\bigr),
\]
which are the paper's two tables, centred, and which lie in $\Theta_{0}$.

\leanbox{\lean{shiftArms}, \lean{rng\_sub\_const}, \lean{energy\_shiftArms},
\lean{inClass\_center}, \lean{thetaLowC\_eq\_centered\_raw},
\lean{low\_world\_hypotheses\_consistent} in \lean{LeanPaperASeparationV1}}

\begin{theorem}[The surrogate is wrong in both directions, at worlds that exist]
\label{thm:sep}
Put variance $M$ on the two cells named by each family and $1$ elsewhere. At $\theta^{-}$,
$\Gamma_{\rm true}/\Gamma_{\rm proxy}\le64/M$; at $\theta^{+}$,
$\Gamma_{\rm true}/\Gamma_{\rm proxy}\ge M/96$. Consequently, for every $c>0$ there is a
single $M$ at which the first ratio falls below $c$ and the second rises above it, so no
constant relates the surrogate to the truth in both families at once.
\end{theorem}

Both games are suprema of infima, so the statement is empty unless they exist. They do: each
inner set is nonempty and bounded below by zero, and each outer set is capped by the cheap
alternative the separation already exhibits, so an explicit bound is available in each of
the four cases.

\leanbox{\lean{low\_separation\_at\_world}, \lean{high\_separation\_at\_world},
\lean{exists\_low\_games}, \lean{exists\_high\_games},
\lean{no\_constant\_factor\_on\_families} in \lean{LeanPaperASeparationV1}}

\NOTEPAR{Scope. $M$ is a known variance of the observation model, so this is a separation
within a heteroscedastic family, at $\varepsilon=0$, at an instance whose decision is
unique. Removing the dependence on heteroscedasticity needs the $r$-pair variant, which is
not proved.}

\paragraph{What is and is not new here.}
That an information proxy can select a design lacking identifiability is
known: Overstall studies expected Fisher-information gain and gives conditions under
which the optimal design induces non-identifiability~\cite{overstall2022}. That an
acquisition rule should target the decision rather than the predictor is
known~\cite{filstroff2021}, as is error amplification under model
misspecification~\cite{tang2026}, and decision-focused sequential design already carries
a sample-complexity guarantee together with an instance on which a prediction-focused
method fails to eliminate a decision-relevant hypothesis~\cite{wan2026}. When a proxy
\emph{suffices} for a decision, and what to do when it does not, is itself a studied
question~\cite{schutte2026}, and targeted active learning for probabilistic models already
carries label-complexity bounds~\cite{tosh2022}. In combinatorial
complexity there is now a total Boolean function with
$R(f)=\widetilde{O}(\sqrt{C(f)})$~\cite{bendavid2026}, which forecloses the generic
headline that certification must cost more than deciding.

So the generic statements are taken. What this paper claims is narrower and, we think,
sharper: the \emph{observation model is correct}, the stopping certificate is
\emph{valid}, and the decision objective is \emph{unchanged}; the only thing altered is
the geometry used to score allocations, and that alone changes what the design can learn.
The failure is not misspecification and not an invalid test --- it is a frozen extremum
in an otherwise correct model.

\subsection{A repair criterion, and exactly what it does not do}

A design carries a certificate when a computable barrier separates the surrogate's value
from the nearest switch,
\[
J_{\rm proxy}(w)<B_{\rm switch}(w)\ \Longrightarrow\ J_{\rm true}(w)=J_{\rm proxy}(w),
\]
which is a theorem. What it is worth is a separate question, and answering it requires
care about the denominator.

Wherever the certificate fires it forces $J_{\rm true}=J_{\rm proxy}$, so a certified
design inherits the separation:
\begin{equation}
\mathrm{Cert}(w)\ \Longrightarrow\ \frac{J_{\rm proxy}(w)}{\Gamma_{\rm proxy}}\;\le\;
\frac{64}{M}.
\label{eq:certceiling}
\end{equation}
It is tempting to read \eqref{eq:certceiling} as saying certified designs become worthless
as $M$ grows. That reading is wrong, and the reason is that $\Gamma_{\rm proxy}$ is the
wrong yardstick: what a repair is for is the \emph{true} game, so the quantity to bound is
$J_{\rm true}(w)/\Gamma_{\rm true}$, and at the optimum level
\[
\frac{\Gamma_{\rm cert}}{\Gamma_{\rm true}},
\qquad
\Gamma_{\rm cert}=\sup_{w\,:\,\mathrm{Cert}(w)}J_{\rm true}(w).
\]
Converting \eqref{eq:certceiling} into those units costs the separation back,
$J_{\rm true}(w)/\Gamma_{\rm true}\le(64/M)\,\Gamma_{\rm proxy}/\Gamma_{\rm true}$, and on
the sweep's own bracket that product is $66.2$, $18.4$, $6.43$, $3.42$ and $2.58$ at
$M=1,4,16,64,256$. It exceeds one at every size: \textbf{in true units the ceiling is
vacuous, and \eqref{eq:certceiling} constrains the repair's efficiency not at all.}
\CERT{}

The same correction applies to the exhibited witness. At $M=256$ the rational design
$w=\tfrac1{1000}(1,1,1,1,1,1,497,497)$ has $J_{\rm proxy}=1/128\,000$ and $B_{\rm
switch}=497/24\,096\,000$, so $J_{\rm proxy}<B_{\rm switch}$ and the criterion
\emph{fires}. That design is $0.402\%$ of the floored surrogate optimum but $4.15\%$ of the
bracketed true optimum --- a factor of ten between the two readings, and the second is the
one that matters. One poor witness is in any case not a statement about
$\Gamma_{\rm cert}$: at $M=256$ none of the $43$ sampled designs is certified, so this
sweep exhibits no certified design other than the witness and gives no lower bound on
$\Gamma_{\rm cert}$ at all.

\NOTEPAR{So the honest position on repair is neither positive nor negative: it is
undetermined. What the criterion gives is a \emph{verified} equality at the designs where
it fires; what is not known is the value of $\Gamma_{\rm cert}/\Gamma_{\rm true}$, because
\eqref{eq:certceiling} is vacuous in true units and the sweep certifies too few designs to
estimate the supremum. Settling it means optimising $J_{\rm true}$ \emph{over the
certificate region}, with $\Gamma_{\rm true}$ as denominator. Do not write that certified
designs are worthless, and do not write that the criterion repairs the surrogate.}

\leanbox{\lean{barrier\_implies\_equality}, \lean{optima\_equal\_of\_certificate},
\lean{certifiable\_fraction\_le\_sixtyfour\_over\_M},
\lean{design\_certificate\_insufficient} in \lean{LeanPaperARepairV1}}

\subsection{A learner that stops under a certificate}
\label{sec:learner}

The separation is about information rates. Turning one into a sample count needs a rule that
actually stops, so we run one: a learner that sees only counts, observations and known
variances, brackets each answer's value by cutting planes at the current estimate, keeps the
answers whose upper bracket survives $\max_{S}l_{S}-2Lr_{t}$, tracks the selected answer's
design with forced exploration, and stops when a \emph{certified lower} bound on every
region's infimum exceeds the threshold. The certificate is weak duality,
$g(y)=-\tfrac12y^{\top}AH^{-1}A^{\top}y+y^{\top}(A\hat\mu-b)\le\min$ for every $y\ge0$, so
it is valid whatever the optimiser does; a primal objective is an \emph{upper} bound and
cannot justify stopping.

On a centred instance with $M=16$, $12$ runs at each of
$\alpha\in\{0.1,0.05,0.02,0.01\}$, under \emph{common random numbers} across the
$\alpha$ panel --- $\alpha$ enters the algorithm only at the stopping test, so a seed free
of $\alpha$ puts the four levels on one path and makes $\tau$ monotone in $\alpha$
pathwise, which the artifact checks on all $36$ paths. All $144$ runs returned the correct
answer with zero regret and none was censored.

\begin{table}[t]
\caption{Mean stopping time per confidence level, $n=12$ per cell, $\pm$ one standard
error. The learner-versus-uniform column is a Welch test with a Holm correction across the
four levels. The pooled mean is deliberately \emph{not} reported as a headline: it averages
four configurations whose individual evidence differs.}
\label{tab:panel}
\centering
\footnotesize
\begin{tabular}{@{}lrrrrr@{}}
\toprule
$\alpha$ & learner & oracle (diag.) & uniform & $t$ & $p$ (Holm) \\
\midrule
$0.10$ & $1316.7\pm88.6$ & $1104.2\pm72.9$ & $1620.8\pm92.4$ & $2.38$ & $0.079$ \\
$0.05$ & $1366.7\pm86.9$ & $1183.3\pm89.9$ & $1650.0\pm91.5$ & $2.25$ & $0.079$ \\
$0.02$ & $1416.7\pm102.7$ & $1212.5\pm93.0$ & $1725.0\pm78.0$ & $2.39$ & $0.079$ \\
$0.01$ & $1429.2\pm102.3$ & $1283.3\pm84.0$ & $1762.5\pm82.1$ & $2.54$ & $0.076$ \\
\bottomrule
\end{tabular}
\end{table}

The learner's mean is below uniform at every level, by $19$--$23\%$, with a strikingly
stable $t\approx2.2$--$2.5$; \emph{no cell reaches significance after the Holm
correction} at $n=12$, and we report it that way rather than quoting the pooled difference.
An earlier, independently-seeded panel showed one cell at $t=4.67$; under common random
numbers that disappears, which is what an unlucky draw looks like.

\paragraph{The $\alpha$ axis as a slope, which is the part that \emph{is} testable.}
The levels cannot be compared with $\log(1/\alpha)/\Gamma$: the threshold carries
$\tfrac12\sum_{e}\log(1+N_{e}/a)$, measured at $19.6$--$19.8$ here against
$\log(1/\alpha)\le4.61$. An additive penalty cancels in the \emph{slope}, and common
random numbers make the within-path difference clean. Regressing $\tau$ on
$\log(1/\alpha)$ along each path gives a learner slope of $49.6\pm12.8$
(95\%\ CI $[24.6,74.7]$). The penalty itself drifts at $+0.113$ per unit of
$\log(1/\alpha)$, so the drift-corrected prediction is $(1+0.113)/\Gamma_{\rm comp}=47.9$
rather than $1/\Gamma_{\rm comp}=43.0$, and the interval covers it. \NUM{}

\NOTEPAR{This is a \emph{heuristic diagnostic}, not a test of the rate, and we say so
because the arithmetic does not support more. Writing
$Z_{\tau}=\Gamma\tau+R_{\tau}$ and
$Z_{\tau}=\log(1/\alpha)+b_{\tau}+O_{\alpha}$ gives
$\Gamma\,\mathbb{E}\tau=\log(1/\alpha)+\mathbb{E}b_{\tau}+\mathbb{E}O_{\alpha}-\mathbb{E}R_{\tau}$.
Differencing along the $\alpha$ axis removes the constant part of the penalty $b_{\tau}$,
and measuring its drift removes the rest of it --- but it removes neither the threshold
overshoot $O_{\alpha}$ nor the finite-time information residual $R_{\tau}$, and neither is
bounded anywhere in this work. Either can move the slope in either direction. So the
agreement above is consistent with the rate theorem and establishes nothing on its own; an
earlier draft called the corrected figure an upper estimate of $1/\Gamma$, which was wrong.
The baseline also spans only $\log(1/\alpha)\in[2.30,4.61]$ and the interval is
$\pm26\%$.}

\paragraph{What the implementation guarantees, and how that is checked.}
Three properties the analysis needs are enforced in the code and checked structurally
rather than asserted. The oracle's own error must vanish for any accumulation argument to
apply, so the cutting-plane tolerance is tightened as $1/n^{2}$ --- floored at a width the
bracket can actually reach, since a tolerance it cannot clear never fires the early break
--- and the iteration budget
raised with the refresh index, and the achieved bracket gap is recorded at every refresh
and checked to decrease --- an iteration cap alone guarantees nothing. Every refresh is
stored and the stored count is checked against the number performed, so no trace is
silently truncated. And the learner cannot read the true world because the world is not
in scope: \texttt{run} receives a sampler closure, \texttt{theta} is not among its
parameters, and the self-check inspects the signature instead of promising it in a
comment. \CERT{}

\NOTEPAR{Read this as what it is. The oracle arm reads the true world and is a
\emph{diagnostic}, not a competitor---and not a lower bound on stopping time either: an
oracle allocation maximises the asymptotic information rate, so at finite $\alpha$ a
learner can stop earlier on a path without anything being wrong. (An earlier draft called
it an ``unattainable lower envelope on stopping time''; that was false and is withdrawn.)
No selection snapshot ever narrowed the answer set---not even on the swapped world whose
correct answer is second in the tie-break order---so these runs are evidence about
\emph{stopping}, not about elimination. And the panel cannot be read as a rate: the
threshold is $\log(1/\alpha)+\tfrac12\sum_{e}\log(1+N_{e}/a)$, whose second term measures
$18.5$--$19.7$ here against $\log(1/\alpha)\le4.61$, so the $\alpha$ axis is far too short
to expose the asymptotic regime. Do not write ``the learner attains $\Gamma_{\rm comp}$''.}

\subsection{The confidence sequence, with Ville proved}
\label{sec:ville}

The stopping rule is correct if the statistic crossing its threshold is unlikely under the
true world, and that reduces to one inequality: a non-negative supermartingale started at
mean one exceeds $1/\delta$ with probability at most $\delta$. Mathlib has Doob's maximal
inequality for non-negative \emph{sub}martingales, which is the wrong direction, and does
not have Ville's. We prove it, from optional stopping:

\begin{theorem}[Ville, at a finite horizon]
\label{thm:ville}
Let $M$ be a non-negative supermartingale. Then for every $\lambda>0$ and every $N$,
\[
\lambda\cdot\Pr\Bigl(\exists\, j\le N:\ M_{j}\ge\lambda\Bigr)\ \le\ \mathbb{E}[M_{0}].
\]
\end{theorem}

The proof is three steps. Negating $M$ gives a submartingale, so optional stopping between
the constant time $0$ and the hitting time of $[\lambda,\infty)$ capped at $N$ gives
$\mathbb{E}[M_{\tau}]\le\mathbb{E}[M_{0}]$. On the event that $M$ reaches $\lambda$ by time
$N$ the stopped value is at least $\lambda$, and off it the value is at least $0$ because
$M$ is non-negative. Combining the two bounds the probability.

The consequence is stated for an actual measure and an actual supermartingale, with no
hypothesis standing in for Ville: if crossing the threshold forces the mixture past
$1/\delta$, the rule errs with probability at most $\delta$.

\leanbox{\lean{ville\_finite\_horizon}, \lean{confidence\_of\_supermartingale},
\lean{stopping\_error\_le} in \lean{LeanPaperA7ConfidenceV1}}

\paragraph{A process it applies to.}
A theorem about supermartingales is worth what its instances are worth, so we exhibit one.
Take $\Omega=\mathbb{N}\to\mathbb{R}$ with the law of independent $\mathcal{N}(0,1)$
observations, let $Y_{n}$ be the $n$-th observation, and let
\[
\Lambda_{n}(u)=\exp\!\bigl(uY_{n}-u^{2}/2\bigr),
\qquad
M_{t}(u)=\prod_{n\le t}\Lambda_{n}(u).
\]
The coordinates are independent because the law is a product; the increments are therefore
independent, being measurable functions of them; and $\mathbb{E}\Lambda_{n}(u)=1$ exactly,
which is the Gaussian moment generating function $\mathbb{E}e^{uY}=e^{u^{2}/2}$ cancelled
against the tilt. Independence makes the conditional expectation of the next increment its
mean, so $M(u)$ is a martingale for the natural filtration --- and Theorem~\ref{thm:ville}
applies to it, giving
$\Pr(\exists\, j\le N:\ M_{j}(u)\ge1/\delta)\le\delta$ on that stream.

\leanbox{\lean{obs\_indep}, \lean{incr\_integral}, \lean{mix\_indep\_incr},
\lean{mix\_martingale\_step}, \lean{mix\_supermartingale}, \lean{stream\_confidence} in
\lean{LeanPaperA7StreamV1}}

\NOTEPAR{Scope of \emph{this} instance. It has a fixed tilt $u$ and queries the same
coordinate at every step, so it is not the learner's process; what it establishes is that the
supermartingale hypothesis is satisfiable by a real observation model and that the Gaussian
step calculation is the right one.}

\paragraph{The process the code actually thresholds, at a random past.}
The learner's statistic is the closed form
$M_{t}=\prod_{e}\sqrt{a/(a+N_{e}(t))}\exp\bigl(S_{e}(t)^{2}/2(a+N_{e}(t))\bigr)$, and the
one-step identity $\mathbb E[f_{n+1}(s+Z)]=f_{n}(s)$ holds at a \emph{fixed} $(n,s)$. The
process needs it at $n=N_{e}(t)$ and $s=S_{e}(t)$, which are measurable in the past. The
tempting shortcut --- that $S_{e}$ given $N_{e}$ is Gaussian --- is \emph{false} under
adaptive sampling, because the count's evolution depends on the observations.

What licenses the substitution is the freezing lemma,
\[
\mathbb E\bigl[f(X,Y)\mid\sigma(X)\bigr]
=\Bigl(x\mapsto\int f(x,y)\,d\mathcal L(Y)\Bigr)\Big|_{x=X},
\qquad Y\perp X,
\]
which Mathlib does not have and which we build: independence makes the joint law a product, a
product is a composition with a constant kernel, and the conditional distribution is
characterised by that composition. With the inner integral evaluated by the fixed-past
identity, $\mathbb E[M_{t+1}\mid\mathcal F_{t}]=M_{t}$.

Integrability is \emph{derived}, not assumed: by the product criterion the joint is integrable
exactly when the previous step is, since the inner integrals are finite and their values are
the previous process --- so $\mathbb E M_{t}=1$ and integrability propagate together from
$M_{0}=1$. Finally $M$ \emph{is} the implementation's mixture, crossing $Q_{t}>b_{t}(\delta)$
forces $M_{t}\ge1/\delta$ because the penalty $\tfrac12\sum_{e}\log(1+N_{e}/a)$ cancels on
both sides, and the one-step equalities assemble into a martingale, hence a supermartingale,
to which Theorem~\ref{thm:ville} applies.

\NOTEPAR{What is still owed is an identification, not an analysis: that the code's $q_{t}$,
$N_{e}(t)$ and $S_{e}(t)$ are the model's. The constant $a$ is fixed and the process is at a
fixed true parameter; substituting a time-varying estimate preserves none of this. The
per-step confidence guarantee at the algorithm's own confidence set also remains.}

\leanbox{\lean{condDistrib\_const\_of\_indepFun}, \lean{condExp\_freeze},
\lean{condExp\_freeze\_apply} in \lean{LeanPaperAFreezeV1};
\lean{martingale\_step}, \lean{mean\_step}, \lean{joint\_integrable},
\lean{inner\_integrable\_law} in \lean{LeanPaperAAdaptMartV1};
\lean{Mval\_eq\_Mstat}, \lean{threshold\_crossing}, \lean{closed\_form\_error\_bound} in
\lean{LeanPaperAStopJoinV1}}

\subsection{Separation without heteroscedasticity: the $r$-pair family}
\label{sec:rpair}

Theorem~\ref{thm:sep} needs $M$ to grow, so at $M=1$ it says nothing and the separation can
be read as an artefact of a badly conditioned observation model. It is not. Widen the
relation instead of the variance: give each arm $2r+2$ actions, with $r$ disjoint pairs of
middle cells in arm $B$. The frozen proxy still watches only the four anchors, while the
adversary may move whichever pair the design has measured least.

\begin{theorem}[The $r$-pair separation]
\label{thm:rpair}
With $r$ disjoint pairs and a unit design budget, some pair carries weight at most $1/r$, so
$\Gamma_{\rm true}\le1/8Mr$ while $\Gamma_{\rm proxy}\ge1/512$. Hence
$\Gamma_{\rm true}/\Gamma_{\rm proxy}\le64/rM$, and for every $c>0$ there is an $r$ at which
the ratio falls below $c$ \emph{at $M=1$}.
\end{theorem}

The pigeonhole is the whole mechanism, and it is what makes the bound hold after the
allocation has been optimised: the adversary reads $w$ and then chooses the pair.

\leanbox{\lean{exists\_light\_pair}, \lean{pair\_shift\_energy},
\lean{light\_pair\_energy\_le}, \lean{rpair\_ratio\_le},
\lean{no\_constant\_factor\_at\_unit\_variance} in \lean{LeanPaperARPairV1}}

\paragraph{The matching lower bound.}
Both sides are now proved, so the ratio is $\Theta\bigl((rM)^{-1}\bigr)$ and not merely
$O$. The lower side needs an explicit design rather than a pigeonhole: the anchor design
spreads mass over the four contrast cells and the $2r$ pair cells, and the four-square
inequality
\[
\tfrac{1}{4}<x_{1}-x_{2}+x_{3}-x_{4}
\;\Longrightarrow\;
x_{1}^{2}+x_{2}^{2}+x_{3}^{2}+x_{4}^{2}>\tfrac{1}{64}
\]
turns a range gap at the contrast into an energy floor. Evaluating $\Gamma_{\rm true}$ at
that design gives
\[
\frac{1}{128\,(4+2rM)}\;\le\;\Gamma_{\rm true}\;\le\;\frac{1}{8rM},
\qquad
\frac{2}{2+rM}\;\le\;\frac{\Gamma_{\rm true}}{\Gamma_{\rm proxy}}\;\le\;\frac{64}{rM}.
\]

\leanbox{\lean{anchorDesign}, \lean{four\_sq\_lower}, \lean{rpair\_true\_ge},
\lean{rpair\_ratio\_ge}, \lean{rpair\_theta} in \lean{LeanPaperARPairV1}}

\paragraph{The two games, as games.}
The bracket above is arithmetic about two numbers: it takes $\Gamma_{\rm true}$ and
$\Gamma_{\rm proxy}$ as given and pins their ratio. That is not yet a statement about the
$r$-pair problem, so we state it as one. Fix the world
$\theta_{0}$ with arm $A$'s anchors a unit apart and all of arm $B$ at zero, and the
variance profile unit on the anchors and $M$ on the pairs. Define
\[
\Gamma_{\rm true}
=\sup_{w\ \rm alloc}\ \inf_{\lambda:\ \mathrm{range}_{A}\lambda<\mathrm{range}_{B}\lambda}
  \mathcal{E}(w,\lambda),
\qquad
\Gamma_{\rm proxy}
=\sup_{w\ \rm alloc}\ \inf_{\lambda\in\Lambda_{\rm proxy}}\mathcal{E}(w,\lambda),
\]
where $\Lambda_{\rm proxy}$ is the set of alternatives the frozen surrogate can separate
from $\theta_{0}$ at all --- those that move an anchor and leave the pair cells alone.
Both are genuine $\sup\inf$ over the actual allocations and the actual alternatives.

\begin{theorem}[The ratio, on the games themselves]
\label{thm:rpairgames}
For $r\ge1$ and $M>0$,
\[
\frac{1}{128\,(4+2rM)}
\;\le\;\frac{\Gamma_{\rm true}}{\Gamma_{\rm proxy}}\;\le\;
\frac{16}{Mr},
\]
and all four one-sided bounds --- $\Gamma_{\rm true}\le1/(8Mr)$,
$\Gamma_{\rm true}\ge1/(128(4+2rM))$, $\Gamma_{\rm proxy}\le1$,
$\Gamma_{\rm proxy}\ge1/128$ --- are theorems, not inputs.
\end{theorem}

The proxy value stays between two absolute constants while the true value falls like
$1/(rM)$: the surrogate overstates the information available by a factor growing without
bound in $r$, at uniform variance $M=1$.

\paragraph{The same bracket, against the surrogate itself.}
$\Lambda_{\rm proxy}$ above pins the middle cells and then uses the \emph{true} ranges;
the frozen-linear surrogate pins the extremum \emph{indices} and tests one linear
contrast, $(\lambda_{1}-\lambda_{0})-(\lambda_{2}-\lambda_{3})\le0$. These are
\emph{not} the same set, and we exhibit the difference rather than wave at it: the world
with arm $A$ at $(0,1)$, arm $B$'s frozen maximum at $-2$ and the middle cells at $0$ has
true $B$-range $2>1$, so it is a restricted-coordinate alternative, while its frozen
contrast is $(1-0)-(-2-0)=3>0$, so it is not a surrogate alternative. Set equality is
false and is not assumed.

What the separation needs is not equality of the sets but the two \emph{values}, and those
we bound directly. Write $\Gamma_{\rm frozen}=\sup_{w}\inf_{\lambda\in\Lambda_{\rm
frozen}}\mathcal{E}(w,\lambda)$.

\begin{theorem}[The bracket holds against the frozen surrogate]
\label{thm:frozengames}
$1/32\le\Gamma_{\rm frozen}\le1$, and hence
\[
\frac{1}{128\,(4+2rM)}\;\le\;\frac{\Gamma_{\rm true}}{\Gamma_{\rm frozen}}\;\le\;
\frac{16}{Mr}.
\]
\end{theorem}

The lower bound is the same mechanism as before: flipping the frozen contrast moves an
alternating sum of four unit-variance cells by one, which by Cauchy--Schwarz costs squared
displacement at least $1/4$, so the uniform-anchor design already secures $1/32$. The
upper bound bumps whichever of arm $B$'s two frozen anchors carries less weight. So the
conclusion is about the surrogate the paper criticises, not a stand-in for it.

\leanbox{\lean{proxyAlt\_ne\_frozenAlt}, \lean{four\_sq\_lower\_one},
\lean{energy\_wUnif}, \lean{frozenValue\_unif\_ge}, \lean{frozenValue\_le},
\lean{frozenGame\_ge}, \lean{frozenGame\_le}, \lean{frozen\_games\_theta} in
\lean{LeanPaperAFrozenV1}}

\paragraph{The world is in the class where the energy is the information.}
`energy` coincides with the Gaussian KL only on the centred class --- each arm's table
summing to zero --- because otherwise the response carries the other arm's mean and that
term enters the divergence. The base world is therefore taken with arm $A$'s anchors at
$(-\tfrac12,+\tfrac12)$ rather than $(0,1)$: this centres both arms while preserving the
unit gap $\theta_{1}-\theta_{0}=1$ that every bound above uses, and it leaves arm $B$ flat
at zero. \lean{theta0\_centered} records the membership, so
\lean{centered\_additive\_response} applies and the games optimise information rather
than a stand-in for it.

\NOTEPAR{What this does \emph{not} do. Centring makes the energy-to-KL identification
available; it does not by itself carry the $r$-pair alternatives through the actual
response map. The alternative sets here are still defined by ranges and by the frozen
contrast, and the step from those to the observation model's own alternative class is not
written.}

\leanbox{\lean{trueGame\_le}, \lean{trueGame\_ge}, \lean{proxyGame\_le},
\lean{proxyGame\_ge}, \lean{rpair\_games\_theta}, with \lean{pairShift\_mem\_alt} and
\lean{anchorBump\_mem} for non-vacuity, in \lean{LeanPaperARGamesV1}}

\NOTEPAR{The two constants, $2$ and $64$, are not matched --- the bracket is
$\Theta((rM)^{-1})$ in order, not a single constant. Do not report a sharp constant.}

\paragraph{From starvation to a leakage bound, without a monotonicity modulus.}
The starvation theorems assume \lean{AnchorMonotone}: raising \emph{any one} anchor weight
strictly increases the surrogate. The real frozen-linear surrogate does not satisfy that
at the boundary --- with another essential anchor still unmeasured, raising the first can
leave the value at zero --- so a bound derived from it would not apply to the object we
are criticising. Positive homogeneity does apply, and suffices.

Let $a=\sum_{\text{anchors}}w_{e}$ and let $\widetilde w$ zero the pair cells and rescale
the anchors by $1/a$. The surrogate reads only the anchors and is homogeneous of degree
one in them, so $J_{\rm proxy}(\widetilde w)=J_{\rm proxy}(w)/a$. If $w$ is within
$\varepsilon$ of the surrogate's optimum then $J_{\rm proxy}(w)/a\le J_{\rm proxy}(w)+
\varepsilon$, and since $a\le1$,
\[
\Bigl(\sum_{\text{middle}}w_{e}\Bigr)\,J_{\rm proxy}(w)\;\le\;\varepsilon,
\qquad\text{hence}\qquad
\sum_{\text{middle}}w_{e}\;\le\;\frac{\varepsilon}{\Gamma_{\rm proxy}-\varepsilon}.
\]
No modulus is assumed: homogeneity is a property of the formula. This is the bridge from
\emph{static} starvation to the \emph{exploration cost} of a policy, because it is the
off-anchor mass, not the design's optimality, that the true alternatives can be measured
with.

\leanbox{\lean{anchorMass\_add\_pairMass}, \lean{renormAnchors\_isAlloc},
\lean{middle\_mass\_le\_of\_homogeneous}, \lean{middle\_mass\_le\_of\_near\_optimum} in
\lean{LeanPaperAHomogV1}}

\subsection{The surrogate misdesigns, not merely misprices}
\label{sec:alloc}

A ratio of optimal \emph{values} can diverge while both objectives still select the same
design: an objective off by a constant picks the right optimiser. So
Section~\ref{sec:rpair} does not yet show the surrogate is harmful to use. The quantity
that does is the efficiency of the design it selects, measured in the true game,
\[
\mathrm{Eff}_{\rm proxy}(\theta)
=\frac{J_{\rm true}(w^{\star}_{\rm proxy};\theta)}{\Gamma_{\rm true}(\theta)} .
\]

On the $r$-pair family this is zero, and for \emph{every} optimiser rather than one
tie-break. Two properties characterise the frozen surrogate without naming its formula:
its value depends only on the four anchor weights, and it strictly prefers anchor mass.
From those alone, any allocation giving a middle pair positive weight can be improved by
sweeping that weight onto an anchor --- the result is still an allocation and the
surrogate strictly increases --- so every point of $\arg\max J_{\rm proxy}$ gives the
pairs weight zero. A design that starves the pairs has \emph{exactly} zero information
against a pair alternative, since that alternative moves only cells the design never
measures:
\[
\forall\,w\in\arg\max J_{\rm proxy},\quad\forall j,d:\qquad
\mathrm{energy}\bigl(w;\theta,\mathrm{pairShift}(\theta,j,d)\bigr)=0
\;<\;\Gamma_{\rm true}.
\]

\leanbox{\lean{proxy\_optimum\_starves\_pairs}, \lean{starved\_energy\_zero},
\lean{proxy\_efficiency\_zero}, \lean{proxy\_true\_value\_zero} in
\lean{LeanPaperARPairV1}}

That argument needs the surrogate to \emph{strictly} prefer anchor mass. On the centred
construction the same conclusion follows from an upper bound alone, with no monotonicity
modulus and no homogeneity hypothesis, and the frozen game's value is exact.

\begin{theorem}[The separation, both sides, on one construction and one class]
\label{thm:sepcentred}
Let $\theta_{C}$ place arm $A$'s extrema at $(-\tfrac12,+\tfrac12)$, arm $B$'s at
$(-\tfrac18,+\tfrac18)$ and its $2r$ middle cells at $0$; let the four extremum cells have
variance $1$ and the middle cells variance $M>0$; and let both games take their infimum over
subsets of the \emph{same} class
$\Theta_{0}=\{\lambda:\ \text{each arm's cells sum to zero}\}$. Then for every $r\ge1$
and every $M>0$,
\[
\Gamma_{\rm frozen}=\frac{9}{512},
\qquad
\frac{9}{128(4+2rM)}\;\le\;\Gamma_{\rm true}\;\le\;\frac{1}{8Mr},
\]
and consequently
\[
\frac{4}{4+2rM}
\;\le\;
\frac{\Gamma_{\rm true}}{\Gamma_{\rm frozen}}
\;\le\;
\frac{64}{9rM} .
\]
\end{theorem}

The gap is $\Theta((rM)^{-1})$ and is present already at $M=1$: it is bought by adding
actions, not by inflating a variance. The upper bound is the light pair --- some pair of
middle cells carries at most $1/r$ of any design, and moving it apart by $\tfrac12$ already
reverses the decision. The lower bound is the variance-proportional design
$w_{e}=\sigma_{e}^{2}/(4+2rM)$, at which the energy is the plain squared distance; reversing
the decision forces a contrast of at least $g=\tfrac34$ across four cells, and
Cauchy--Schwarz turns that into $9/64$ of squared distance. The same $g$ sets
$\Gamma_{\rm frozen}=g^{2}/32$, so one number drives both sides.

\NOTEPAR{Two things this theorem is careful about, because both are easy to get wrong.
First, \emph{both} infima are restricted to $\Theta_{0}$. Restricting the set an infimum
ranges over can only \emph{raise} it, so a bound proved over the unrestricted set does not
transfer; it transfers here because the witnesses --- the contrast-closing shift and the
pair shift --- are themselves centred, which is proved, not observed. Second, the
alternative sets are taken closed. That is what Theorem~\ref{thm:closure} (\lean{strict\_inf\_eq\_closed\_inf}) licenses, and it is
not doing the work: for every $\delta>0$ the shift by $\tfrac12+\delta$ is a
\emph{strictly} alternative world lying in $\Theta_{0}$, so the bound is approached from
inside the strict set.}

\leanbox{\lean{frozenGameC₀\_eq}, \lean{trueGameC₀\_le}, \lean{wCert\_isAlloc},
\lean{energy\_wCert}, \lean{contrast\_movesC}, \lean{trueValueC₀\_wCert\_ge},
\lean{trueGameC₀\_ge}, \lean{centred\_separation}, \lean{pairShiftC\_strict\_alt},
\lean{frozenValueC\_le\_restricted} in \lean{LeanPaperACentredV1}}

\begin{theorem}[The surrogate's optimum starves the cells that carry the compression]
\label{thm:starve}
Let $\theta_{C}$ place arm $A$'s extrema at $(-\tfrac12,+\tfrac12)$, arm $B$'s at
$(-\tfrac18,+\tfrac18)$ and every middle cell at $0$, so both arms are centred and the
frozen contrast has margin $g=\mathrm{range}(A)-\mathrm{range}(B)=\tfrac34$. Then for every
$M>0$,
\[
\Gamma_{\rm frozen}=\frac{g^{2}}{32}=\frac{9}{512},
\]
attained at the uniform-anchor design; and for every allocation $w$,
$J_{\rm frozen}(w)\le\tfrac{9}{512}\,a(w)$ with $a(w)$ the mass on the four anchors.
Consequently
\begin{enumerate}
\item[(i)] every $w$ with $J_{\rm frozen}(w)\ge\tfrac{9}{512}$ has $a(w)=1$ and puts
\emph{exactly zero} mass on the $2r$ middle cells; and
\item[(ii)] every $w$ within $\varepsilon$ of the frozen optimum leaks at most
$\varepsilon/\Gamma_{\rm frozen}=\tfrac{512}{9}\varepsilon$ onto them.
\end{enumerate}
\end{theorem}

Part~(ii) is the link the sequential argument needs: an \emph{approximate} optimiser of the
surrogate---which is what a learner running on the surrogate actually has---still starves
the informative cells, at a rate controlled by its own suboptimality. The starved cells are
the only ones a pair alternative moves, so the true information against that alternative is
zero at (i) and $O(\varepsilon)$ at (ii).

\leanbox{\lean{thetaC\_centred}, \lean{frozen\_margin}, \lean{energy\_lamStar},
\lean{frozenValueC\_le}, \lean{frozenGameC\_eq}, \lean{starvation\_at\_optimum},
\lean{near\_optimal\_starvation} in \lean{LeanPaperACentredV1}}

\NOTEPAR{$J_{\rm true}(w^{\star}_{\rm proxy})=0$ does \emph{not} by itself give
$\mathbb{E}\tau=\Omega(\log^{2}(1/\alpha))$, and we do not claim it. A learner may add
forced exploration, its estimated extrema may move, or an approximate oracle may put mass
on the starved cells; the rate depends on those. The honest target is a law in the
\emph{cumulative exploration}: change of measure gives
$\mathbb{E}N_{\rm movable}(\tau)\ge(2r/d^{2})\,\mathrm{kl}(1-\alpha,\alpha)$, and IF the
sampling rule satisfies $\mathbb{E}N_{\rm movable}(\tau)\le C_{r}+C'_{r}\mathbb{E}[\tau^{a}]$
with $0<a<1$, Jensen gives
$\mathbb{E}\tau\ge[(2r\,\mathrm{kl}(1-\alpha,\alpha)/d^{2}-C_{r})/C'_{r}]_{+}^{1/a}$ ---
so $O(\sqrt t)$ exploration yields $\Omega(\log^{2}(1/\alpha))$ and linear exploration
yields only $\Omega(\log(1/\alpha))$. The hard half is the count bound, not Jensen. It is
proved in Section~\ref{sec:adaptive}, both for the learner that tracks an exact surrogate
optimum fixed at the true world and for the learner that re-solves the surrogate at its
estimate. That proof does not use the chain here: it bounds the middle count path by path.}

\subsection{The alternative set is exactly a union of polyhedra}
\label{sec:region}

The oracle solves the inner problem region by region. That is an oracle for
\emph{compression} only if the decomposition is an equality with the set of worlds at which
the retained answer is wrong --- both inclusions, and derived from the loss rather than
posited. Theorem~\ref{thm:modelbridge} supplies the missing half.

\begin{theorem}[The decomposition, as an equality]
\label{thm:region}
For every tolerance $\varepsilon$ and every $\lambda\in\Theta_{0}$, retaining $A$ is
$\varepsilon$-unsafe if and only if there are arm-$B$ cells $p,q$ with
\[
\lambda_{1}-\lambda_{0}+2\varepsilon<\lambda_{p}-\lambda_{q}
\quad\text{and}\quad
\lambda_{0}-\lambda_{1}+2\varepsilon<\lambda_{p}-\lambda_{q}.
\]
Each pair therefore contributes one polyhedron cut out by two linear inequalities, and the
unsafe set is exactly their union.
\end{theorem}

Forward, take the cells realising arm $B$'s supremum and infimum, whose difference is the
range; backward, any pair's gap is at most the range. The split into two linear inequalities
is $|x|<y\iff x<y\wedge-x<y$.

\NOTEPAR{The absolute value is where such a cover goes wrong, and we report this because it
happened to us. Building the numerical certificate for this same set, we first covered it
with one half-space per sign --- the two inequalities as \emph{alternatives}. That cover
contains the true world itself: take $p,q$ two middle cells, both at $0$, and the sign
pointing away from arm $A$'s gap. The certificate collapses to zero and the rule never
stops. The constraints are simultaneous, not alternative, and the repaired certificate
reproduces the lower bound of Theorem~\ref{thm:sepcentred} exactly at the certifying
design.}

\leanbox{\lean{range\_gap\_iff\_exists\_pair}, \lean{pair\_region\_linear},
\lean{unsafe\_iff\_union\_of\_polyhedra}, \lean{polyhedron\_nonvacuous} in
\lean{LeanPaperARegionV1}}

\paragraph{Two traps in the region decomposition.}
Turning each region into a convex QP needs two steps that are routinely got wrong, and both
change the answer. First, replacing $>$ by $\ge$ is not free: a region's closure may contain
points its strict part does not. What licenses it is a segment. If the region has a strictly
bad $\lambda^{\circ}$ then, for a boundary point $x$, the segment
$(1-t)x+t\lambda^{\circ}$ stays in the ordering polyhedron by convexity, is strictly bad for
every $t>0$ because the badness constraint is affine, and converges to $x$ by continuity ---
so a lower bound valid on the strict region is valid on the closure and the two infima agree.
A region whose strict part is \emph{empty} certifies nothing and must be dropped. Second, a
coordinate with $w_{e}=0$ does not appear in the objective at all, so the inner problem is
unconstrained in it and any inverse-Hessian formula is inapplicable there.

\leanbox{\lean{exists\_strict\_near}, \lean{lower\_bound\_transfers},
\lean{empty\_strict\_contributes\_nothing}, \lean{objective\_indep\_of\_zero\_weight},
\lean{zero\_weight\_unbounded\_below\_direction} in \lean{LeanPaperAOracleRegionV1}}

\subsection{The repair frontier: what buying back information costs}
\label{sec:frontier}

Theorem~\ref{thm:starve} is negative. The operational question is the exchange rate: if one
gives up a fraction $\beta$ of the surrogate's value --- the quantity a practitioner
controls --- how much true information can be recovered? Write
\[
\mathcal F(\beta)=\sup\Bigl\{\,J_{\rm true}(w)\ :\ w\in\Delta,\
J_{\rm fr}(w)\ge(1-\beta)\Gamma_{\rm fr}\,\Bigr\}.
\]

One structural fact settles both sides. Each game's value is an infimum of functions
\emph{linear} in the design, hence \emph{concave} in the design, so mixing two designs
cannot lose more than linearly.

\begin{theorem}[The frontier, pinned from both sides at one order]
\label{thm:frontier}
For every $r\ge1$, $M>0$ and $\beta\in[0,1]$,
\[
\frac{9\beta}{128(4+2rM)}
\;\le\;
\mathcal F(\beta)
\;\le\;
\frac{\beta}{8rM},
\qquad\text{so}\qquad
\mathcal F(\beta)=\Theta\!\left(\frac{\beta}{rM}\right)
\ \text{ for } rM\ge1 .
\]
The lower bound is attained by the explicit mixture
$w_{\beta}=(1-\beta)\,w_{\rm unif}+\beta\,w_{\rm cert}$ of the surrogate's optimum with
the design certifying the true game's lower bound.
\end{theorem}

The upper side is Theorem~\ref{thm:starve}(ii) with $\varepsilon=\beta\Gamma_{\rm fr}$. The
lower side is concavity twice: $J_{\rm fr}(w_{\beta})\ge(1-\beta)\Gamma_{\rm fr}$ makes
$w_{\beta}$ feasible, and $J_{\rm true}(w_{\beta})\ge\beta J_{\rm true}(w_{\rm cert})$
gives the value. At $\beta=0$ the upper bound is $0$, so the frontier \emph{collapses to}
Theorem~\ref{thm:starve}: the two statements are one family.

\NOTEPAR{What is and is not new here. Convex-combining two designs is routine and we do not
claim it. The content is that on this problem the trade-off is pinned from \emph{both} sides
at the \emph{same} order, and that the order is the separation factor itself: the price of
repair is linear in the sacrifice and carries the same $1/(rM)$ as the mispricing. The
statement is about information at a design; it is \emph{not} a sample-complexity claim.}

\leanbox{\lean{mixW\_isAlloc}, \lean{energy\_mixW}, \lean{value\_concave},
\lean{frontier\_upper}, \lean{frontier\_lower}, \lean{frontier\_two\_sided},
\lean{frontier\_at\_zero} in \lean{LeanPaperAFrontierV1}}

\subsection{The construction is an observation model, not only a parameter class}
\label{sec:model-bridge}

The two games are infima of an energy over subsets of a parameter class. For that to be a
statement about \emph{compression}, the class, the response and the observation law have to
be the ones the compression loss is written with. Exhibiting the true world inside the class
does not establish it.

The response is additive, $F_{\lambda}(a,b)=\lambda_{A}(a)+\lambda_{B}(b)$, against the
uniform product reference. Retaining an arm means explaining the response by that arm's own
marginal and paying the oscillation of what is left.

\begin{theorem}[The class, the marginals, the loss and the decision]
\label{thm:modelbridge}
$\lambda\in\Theta_{0}$ \emph{if and only if} both marginals recover their coordinates,
$Q^{\lambda}_{A}=\lambda_{A}$ and $Q^{\lambda}_{B}=\lambda_{B}$. On $\Theta_{0}$ the
residual after retaining $A$ is $\lambda_{B}$ itself, so
\[
L_{\lambda}(\{A\})=\tfrac12\mathrm{range}(\lambda_{B}),
\qquad
L_{\lambda}(\{B\})=\tfrac12\mathrm{range}(\lambda_{A}),
\]
and retaining $A$ is $\varepsilon$-unsafe if and only if
$\mathrm{range}(\lambda_{B})>\mathrm{range}(\lambda_{A})+2\varepsilon$.
\end{theorem}

Centring is therefore not a convenience attached to the construction: it is exactly the
condition under which the uniform-product marginal \emph{is} the coordinate, which is why
the equivalence is the right form for it.

\NOTEPAR{The sign, and what it leaves open. The unsafe condition is \emph{strict}; at
$\varepsilon=0$ it reads $\mathrm{range}(\lambda_{A})<\mathrm{range}(\lambda_{B})$. The
closed set $\{\mathrm{range}(\lambda_{A})\le\mathrm{range}(\lambda_{B})\}$, over which
Theorems~\ref{thm:sepcentred} and~\ref{thm:frontier} are solved, therefore \emph{strictly
contains} the unsafe set: a tie lies in it and is not unsafe. That the two infima agree is
what Theorem~\ref{thm:closure} has to license, and it is \emph{not} yet proved here;
\lean{pairShiftC\_strict\_alt} supplies strictly unsafe witnesses inside $\Theta_{0}$,
which is the ingredient and not the conclusion. The identification of the energy with the
Gaussian divergence for \emph{every} $\lambda\in\Theta_{0}$ is likewise still owed.}

\leanbox{\lean{centred\_iff\_marginals}, \lean{residA\_eq}, \lean{rangeB\_eq\_osc},
\lean{lossRetainA\_eq}, \lean{unsafeA\_iff\_range\_gap}, \lean{unsafeA\_zero\_iff},
\lean{unsafe\_subset\_closed} in \lean{LeanPaperACentredModelV1}}

\subsection{The adaptive query, inside the conditional expectation}
\label{sec:adaptive}

The martingale of Section~\ref{sec:ville} queries one coordinate. The learner picks $q_{t}$
from the history, and the step that must survive is the conditional expectation. It does,
by decomposing over the finitely many queries:
\[
W_{t}(u)=\sum_{e}\mathbf{1}\{q_{t}=e\}\,
\exp\!\bigl(u_{e}Y_{t}-u_{e}^{2}/2\bigr).
\]
The indicator is measurable with respect to the past, so it leaves the conditional
expectation; the exponential is a function of the fresh innovation, so it is independent of
the past and its conditional expectation is its mean, which is one. What is left is
$\sum_{e}\mathbf{1}\{q_{t}=e\}=1$.

The product $L_{t}=\prod_{s\le t}W_{s}$ is then a martingale, with integrability
\emph{derived} rather than assumed --- the past factor is $\mathcal{F}_{t}$-measurable, each
fixed-tilt increment is a function of the fresh innovation, and independence of the two
sigma-algebras descends to the two functions. Ville applies, giving the anytime bound on an
adaptively queried stream.

Three hypotheses carry the policy, and together they are what ``adapted and predictable''
means: the increment at time $s$ is determined by the innovations up to $s$; the query at
$s+1$ by those up to $s$; and the first query is constant, because it is chosen before
anything is seen.

\leanbox{\lean{condExp\_adaptIncr}, \lean{indepFun\_past\_innov},
\lean{adaptMix\_supermartingale}, \lean{adaptive\_confidence} in
\lean{LeanPaperA7StreamV1}}

\paragraph{The prior, and why a fixed grid is the whole obstruction.}
Mixing over $u$ on a finite grid is legitimate --- a finite mixture of supermartingales is
one --- but a grid whose mesh does \emph{not} move misses the exponent, and the failure is
sharp: if $d$ is the smallest non-zero tilt and the true gap satisfies $0<\Delta<d/2$, then
every non-zero grid tilt has rate $u\Delta-u^{2}/2<0$ while the true information is
$\Delta^{2}/2>0$.

What repairs it is the identity
\[
u\Delta-\tfrac{1}{2}u^{2}\;=\;\tfrac{1}{2}\Delta^{2}-\tfrac{1}{2}(u-\Delta)^{2},
\]
which says the rate is the information minus the \emph{squared} miss. A grid of mesh $m$
therefore loses only $m^{2}/2$, second order, so a mesh shrinking with the level at any
rate recovers the exponent; and the price, $\log$ of the number of grid points, is
$\log(2BL+1)=O(\log L)$ against $L=\log(1/\alpha)$, hence negligible. Both halves are
proved, so the fixed mesh is the entire obstruction and conditional Fubini is not needed.
Figure~\ref{fig:a7chain} shows both halves, and the window tail the next paragraph
needs. \NUM{}

\leanbox{\lean{fixed\_grid\_rate\_negative}, \lean{tiltRate\_eq},
\lean{mesh\_grid\_attains\_rate}, \lean{grid\_penalty\_negligible} in
\lean{LeanPaperA7StreamV1}}

\paragraph{From coverage to the expected stopping time.}
The wiring is now a complete deterministic chain. On a window $[n,t)$ on which the
confidence set holds at every step, the accumulated statistic satisfies
\[
\frac{Z_{t}}{t}\;\ge\;\Gamma^{*}-3L_{J}\bar r_{n}-\bar\varepsilon_{n}
-\Gamma^{*}\bar\rho_{n}-\frac{2\Gamma^{*}n}{t}-\frac{C}{t},
\]
with $L_{J}$ the Lipschitz constant of $J$ in the world (not the $L$ of the previous
paragraph), $\bar r_{n},\bar\varepsilon_{n},\bar\rho_{n}$ the window's worst confidence
radius, selection slack and tracking discrepancy, and $C=L_{w}C_{Q}+C_{0}$ the
initialisation loss. Not having stopped by $t$ therefore forces one of the window's
confidence events to have failed --- the union bound's hypothesis is \emph{derived}, not
assumed, which is the step that used to be carried as one. With the window at
$n=\lceil\sqrt t\,\rceil$ and a per-step failure mass $s^{-4}$, the tail is
$\Pr(\tau>t)\le 8t^{-3/2}$. That forces $\Pr(\tau=\infty)=0$, and summing it by its own
telescoping bounds every truncated expectation by $\log(1/\alpha)/\Gamma^{*}+19$; since
$\tau$ is then finite almost surely, monotone convergence lifts the same bound to
$\mathbb{E}\tau$, and hence
$\limsup_{\alpha\to0}\mathbb{E}\tau/\log(1/\alpha)\le1/\Gamma^{*}$.

\leanbox{\lean{good\_window\_rate}, \lean{stopping\_subset\_window\_failure},
\lean{window\_union\_bound}, \lean{sqrt\_window\_tail},
\lean{tail\_forces\_finite\_stopping}, \lean{sqrt\_tail\_sum\_le},
\lean{expected\_stopping\_bound\_at\_level}, \lean{a7\_rate\_limsup} in
\lean{LeanPaperA7TailV1}}

\paragraph{The cost law in cumulative exploration.}
$J_{\rm true}(w^{\star}_{\rm proxy})=0$ does \emph{not} by itself give
$\mathbb{E}\tau=\Omega(\log^{2}(1/\alpha))$: the algorithm may add linear exploration, the
estimated extrema may move, and an approximate oracle may hand mass to cells the surrogate
ignores. The rate turns on exactly those, so the statement is a law in the
\emph{cumulative exploration} and $\log^{2}$ is a corollary of it at one exponent.

Four links. Two were already available: a neighbourhood of the true world on which every
proxy optimum has the same extrema and the same starved support, and the accounting that
telescopes the movable-cell count into $\#\{\text{bad rounds}\}+\sum_{s}(\rho_{s}+
\varepsilon_{s})$. The other two are supplied here. If the per-round probability of
leaving that neighbourhood decays at the inverse-square rate the confidence sequence gives,
the expected number of such rounds is bounded by a constant \emph{uniformly in the
horizon} --- by telescoping $1/(s+2)^{2}\le1/(s+1)-1/(s+2)$. And an $\varepsilon$-optimal
design for the surrogate leaks only $O(\varepsilon)$ off the anchors, by the homogeneity
bound above. Assembling,
\begin{equation}
N_{\rm mov}(t)\;\le\;C_{r}+2\sqrt{t}+\bar\varepsilon\,t .
\label{eq:cumexp}
\end{equation}
The Jensen inversion then runs at $a=1/2$ --- giving $\Omega(\log^{2}(1/\alpha))$ ---
exactly while $\bar\varepsilon t\le2\sqrt t$, that is for $t\le(2/\bar\varepsilon)^{2}$;
past that the linear term takes over and only $\Omega(\log(1/\alpha))$ follows. At the
bracket width measured after the correction below, the crossover sits beyond
$1.5\times10^{13}$ rounds, so on any horizon run here the operating exponent is $1/2$.

\NOTEPAR{What \eqref{eq:cumexp} does not include. (A1), the change-of-measure bound
$\mathbb{E}N_{\rm mov}(\tau)\ge(2r/d^{2})\,\mathrm{kl}(1-\alpha,\alpha)$, is an input. The
inverse-square per-round rate is assumed at the rate the confidence sequence supplies, not
instantiated at the algorithm's own confidence set. So this is the \emph{shape} of the cost
law with its exponent identified, not a finished lower bound on $\mathbb{E}\tau$.}

\leanbox{\lean{inv\_sq\_shift\_sum\_le}, \lean{expected\_bad\_rounds\_bounded},
\lean{offsupport\_mass\_le}, \lean{cumulative\_exploration\_bound},
\lean{sqrt\_dominates\_until}, \lean{measured\_crossover\_is\_far} in
\lean{LeanPaperA2ChainV1}}

\paragraph{What the surrogate learner pays, as an order.}
Starvation is a statement about a design; the operational consequence is an order in
$\log(1/\alpha)$, and it follows from two inputs we name rather than assume away: the
learner's own middle-cell budget $\mathbb E N_{\rm mid}(\tau)\le C+C'\,\mathbb E\sqrt\tau$,
which forced exploration at rate $t^{-1/2}$ produces, and the stopped change of measure at
the hidden pairs. Jensen turns them into
\[
\mathbb E\tau\ \ge\ \Bigl(\frac{2r d^{-2}\,\mathrm{kl}-C}{C'}\Bigr)^{2},
\]
and $\mathrm{kl}\ge(1-\alpha)\log(1/\alpha)-2/e$ makes the right-hand side quadratic in
$\log(1/\alpha)$, against the true-oracle learner's linear cost. The ratio grows without
bound.

\NOTEPAR{This is a statement about a learner with \emph{this} oracle, \emph{this} exploration
rate and \emph{this} tracker. It is not a statement about every method that uses a surrogate,
and the change-of-measure input is exactly what must be supplied for the concrete policy
before it becomes an algorithm-level claim. Both inputs are now supplied, for the learner that tracks an exact surrogate optimum fixed
at the true world and for the learner that re-solves the surrogate at its estimate; see the
last two paragraphs of Section~\ref{sec:adaptive}.}

\leanbox{\lean{surrogate\_cost\_quadratic}, \lean{kl\_at\_correctness},
\lean{surrogate\_omega\_log\_squared}, \lean{separation\_diverges} in
\lean{LeanPaperASurrogateV1}}

\begin{figure}[t]
\centering
\includegraphics[width=\columnwidth]{paperA_swap.pdf}
\caption{Does the rule read the answer off the data? Left: with a certificate for each
answer, every run returns the correct answer on both the base world (answer $A$) and the
swapped world (answer $B$). Right: the rule that can only certify $A$ never stops on the
swapped world --- every run is censored, drawn hatched and labelled rather than placed at the
budget. That is why a panel run only on the base world cannot distinguish learning the answer
from naming it in advance.}
\label{fig:aswap}
\end{figure}

\paragraph{The rate at $\Gamma-\eta$, and what stops it reaching $1/\Gamma$.}
The chain above fixes $t_{0}=\lceil\log(1/\alpha)/\Gamma\rceil+2$, which does not pay for
the code's own $\tfrac12\sum_{e}\log(1+N_{e}/a)$ penalty, nor for estimation, oracle and
tracking error. We therefore prove the $\eta$-parametrised form instead. For every
$0<\eta<\Gamma$: a per-round increment of at least $\Gamma-c\varepsilon$ accumulates at
that rate; the expected-stopping bound holds for \emph{any} non-negative budget $b$, so
the extra threshold terms go inside $b$ rather than being dropped; with
$b=\log(1/\alpha)+s(\alpha)$ and $s(\alpha)=o(\log(1/\alpha))$ the ratio is eventually
within $\epsilon$ of $1/(\Gamma-\eta)$; and $1/(\Gamma-\eta)\to1/\Gamma$ as $\eta\downarrow0$.

What blocks $\eta\downarrow0$ is the oracle's bracket, and we measure it rather than assume
it away --- with one logical caution that an earlier draft missed. A bracket of width
$\varepsilon$ gives $0\le\Gamma-J(w)\le\varepsilon$; it does \emph{not} give
$\Gamma-J(w)\ge\varepsilon$, since the design may be exactly optimal with all the slack in
a loose upper certificate. A positive gap therefore does not show the learner is slower.
What it shows is that the rate we can \emph{certify} degrades, to
$1/(\Gamma-c\varepsilon)$.

The width itself was mis-diagnosed and the correction is recorded rather than quietly
applied. As shipped the bracket settled at $8.0\times10^{-3}$ against
$\Gamma_{\rm comp}\approx2.3\times10^{-2}$, and we read that plateau as the duality gap
between an LP relaxation above and a weak-duality certificate below --- i.e.\ structural.
It is not. Taking the running \emph{minimum} of the relaxation values, damping the step
toward the relaxation's vertex, and solving the regional dual to tolerance give
$1.9\times10^{-4}$ and then $5.0\times10^{-7}$, flat in both the cutting-plane count and
the dual budget. Flatness in both axes rules out those two as the binding constraint; it
does not identify the residual, since the weight floor and the subtracted floor correction
remain candidates and the present checks do not separate them. The honest summary is that
the certified shortfall fell by four orders of magnitude, and that $\eta\downarrow0$
exactly is not proved and cannot be by a floating-point bracket.

\leanbox{\lean{increment\_deficit\_gives\_eta},
\lean{expected\_stopping\_bound\_general}, \lean{rate\_limsup\_eta},
\lean{eta\_to\_zero}, \lean{rate\_floor\_of\_eps\_floor},
\lean{shipped\_floor\_material} in \lean{LeanPaperA7EtaV1}}

\paragraph{The chain, joined --- and the window head it exposes.}
Each link above was proved in isolation and nothing joined them, so the expected-stopping
bound the rate consumes was an \emph{input} to the development rather than an output of it.
Joining them exposes a cost that isolation hid. The inclusion
$\{\tau>t\}\subseteq\bigcup_{s\in[n,t)}\mathrm{good}(s)^{c}$ needs the threshold strictly
below $Gt-2\Gamma n-C$ at the window $n=\lceil\sqrt t\rceil$, so the head of the window must
be paid for: at budget $b$ the crossing time is $b/G+O(\sqrt b)$, not $b/G+O(1)$. We pay it
explicitly rather than absorb it into an asymptotic. For any $\varepsilon\in(0,1)$,
\[
2\Gamma\sqrt t\ \le\ \varepsilon G t+\frac{\Gamma^{2}}{\varepsilon G},
\]
which is $(\varepsilon G\sqrt t-\Gamma)^{2}\ge0$ divided by $\varepsilon G$. Writing
$\mathrm{winBudget}=C+2\Gamma+\Gamma^{2}/(\varepsilon G)+14(1-\varepsilon)G$ --- whose last
summand is what forces the crossing time past $16$, the smallest $t$ at which the $\sqrt t$
window is admissible at all --- the truncated expectations obey
$\mathbb{E}\tau\le(b+\mathrm{winBudget})/((1-\varepsilon)G)+19$, and the rate follows at
$1/((1-\varepsilon)\Gamma)$. What is now a theorem is the COMPOSITION of the links; its
inputs are hypotheses about abstract processes $z$, $Z$ and $\mathrm{good}(s)$, and they are
not yet discharged for the actual learner (see the note below).

\NOTEPAR{$\varepsilon=0$ is not available inside this argument: the split's constant
$\Gamma^{2}/(\varepsilon G)$ blows up. So the joined chain lands in the same $\Gamma-\eta$
family as before, and \lean{chain\_rate\_is\_deficient} records that
$1/\Gamma<1/((1-\varepsilon)\Gamma)$ strictly. Do not read the composition as proving
$1/\Gamma$. The $s^{-4}$ confidence step remains a hypothesis here too.}

\leanbox{\lean{sqrt\_head\_split}, \lean{window\_crossing},
\lean{crossing\_past\_sixteen}, \lean{chain\_tail\_bound},
\lean{expected\_stopping\_bound\_window}, \lean{chain\_expected\_bound},
\lean{a7\_chain\_rate}, \lean{chain\_rate\_is\_deficient} in \lean{LeanPaperA7ChainV1}}

\NOTEPAR{What the chain still takes as input --- all of it, not only the confidence step.
(i) $\Pr(\mathrm{good}(s)^{c})\le s^{-4}$ for the concrete tracking statistic: a Ville
application whose theorem is proved (\lean{ville\_finite\_horizon}) but not instantiated at
the algorithm's own confidence set. (ii) The accumulation hypotheses $G\le z_{s}$ on good
rounds and $Z_{t}\ge\sum_{s<t}z_{s}-C$ for the learner's ACTUAL statistic: this is the
tracking, selection, oracle-perturbation, forced-mixture and concavity chain of the
operational analysis, plus the stopping-certificate gap $\nu_{t}$. Its last link, from
the linear-count energy to the implemented $Q_{t}$ uniformly over the alternatives, is now
proved (\lean{sInf\_Qstat\_ge}, with the $v_{\min}$ bound
\lean{perturbationCost\_div\_le}); the earlier links are not. (iii) The hitting time here is
$\mathbb{N}$-valued; the operational target \lean{family\_operational\_finish} is stated for
the untruncated $\mathbb{N}_{\infty}$ hitting time and still takes the hypotheses
\texttt{hfinite}, \texttt{hUpper} and \texttt{hLower}. So for this abstract chain, and for the SHIPPED learner, the
rate is a CONDITIONAL statement; do not write that it is proved for the shipped learner, and
do not write that only one input remains. The ideal learner of the next paragraph is a
different algorithm, for which all three inputs are discharged.}

\paragraph{The ideal learner, end to end.}
The exact endpoint of the precision contract is a learner whose infima are exact. We define
it completely and prove the rate for it, with no hypothesis beyond $M>0$, $r>0$, $a>0$. After
a warm start of one query per cell, round $s$ selects a pair (answer, design) by a
measurable first maximiser of the exact value $J$ over the answers and a grid of allocations
of mesh $1/(s+1)$; it mixes in the uniform design with weight $\rho_{s}=1/\sqrt{s+1}$ and
tracks the averaged target by a post-warm-up quota; it stops at the first $t\ge K$ at which
some answer's generalised likelihood ratio, the exact infimum of $Q_{t}$ over that answer's
alternatives, exceeds $\log(1/\delta)+\tfrac12\sum_{e}\log(1+N_{e}/a)$. For this learner:
(i) the $s^{-4}$ confidence event is Ville's inequality at the learner's own statistic;
(ii) on that event the statistic accumulates at $\Gamma-\eta$ for every $\eta>0$, with no
box on the worlds; (iii) $\mathbb{E}\tau_{\delta}$ is finite and at most
$\log(1/\delta)/(\Gamma-\eta)+A_{\eta}$. For the lower side the learner is run as a policy of
the kernel protocol that queries a FRESH dummy after stopping. It is coupled path by path
with the unstopped run, so its real-cell count within any horizon is at most $\tau_{\delta}$,
and the policy lower bound applies to it. On the separating family the learner's game is the
actual Gaussian KL game, and
\[
\frac{\mathbb{E}\,\tau_{\delta}}{\log(1/\delta)}\ \longrightarrow\ \frac{1}{\Gamma_{\rm actual}},
\qquad \Pr_{\theta_{C}}(\text{wrong answer})\le\delta .
\]

\NOTEPAR{Scope. This is the IDEAL learner: its oracle maximises the exact $J$ over a refining
net and its stop uses the exact infimum, so the stopping-certificate gap is $\nu_{t}=0$ by
construction. It is not \texttt{EXP\_paperA7\_learner.py} (floor $10^{-6}$, 80-iteration
cap), whose rate remains conditional as stated above. No numerical artifact in this paper
was produced by the ideal learner. Correctness is proved at every world where exactly one
answer is correct (\lean{ideal\_correct}); on the family these are the centred worlds with a
strict range gap, $\theta_{C}$ among them. The finite-expectation and rate statements are
proved at $\theta_{C}$ only.}

\leanbox{\lean{ideal\_family\_finish}, \lean{ideal\_family\_correct}, \lean{GamT\_fam},
\lean{ideal\_family\_lower} in \lean{LeanPaperA7IdealFamilyV1};
\lean{ideal\_expect\_bound}, \lean{ideal\_upper\_rate}, \lean{ideal\_finite},
\lean{ideal\_correct} in \lean{LeanPaperA7IdealUpperV1}; \lean{ideal\_direct} in
\lean{LeanPaperA7IdealDirect2V1}; \lean{runS\_eq}, \lean{cnt\_le\_tau},
\lean{ansEv\_alt\_le}, \lean{ansEv\_true\_ge} in \lean{LeanPaperA7IdealStoppedV1}}

\paragraph{What tracking the surrogate's optimum costs.}
Take any exact optimum $w^{\star}$ of the frozen surrogate game on the family. By the
starvation theorem it puts zero weight on every middle cell, and one exists (the uniform
design on the anchors). Run the learner that tracks it: the same warm start and forcing
$\rho_{s}=1/\sqrt{s+1}$ as the ideal learner, the query
$\arg\max_{e}(s\,w^{\star}_{e}-N_{e})$, and the same GLR stop at level $\delta$ against the
TRUE alternatives. This learner is defined, and its $\tau_{\delta}$ is the hitting time of
that stop. With $K=4+2r$ cells,
\[
\mathbb{E}\,\tau_{\delta}\ \ge\ \Bigl(\frac{KM}{8}\Bigr)^{2}\log^{2}(1/\delta)
\qquad\text{for all }\delta<\delta_{0},
\]
against $O(\log(1/\delta))$ for the ideal learner on the same family, so for every $B$ the
surrogate-tracking cost eventually exceeds $B$ times the ideal cost. The proof takes three
steps, each proved. First, the world that moves one middle pair apart by $1$ is a strict
alternative, and only that pair's queries separate it from $\theta_{C}$, each carrying
divergence $1/(2M)$. Second, forcing is the only source of those queries, so along every
path the pair's count up to $\min(n,\tau_{\delta})$ is at most
$4+2(\tau_{\delta}/s+s)/K$ for every $s>0$; the stopped run's counts equal the unstopped
run's at $\min(n,\tau_{\delta})$, cell by cell. Third, the stopped change of measure at level
$2\delta$, with the horizon chosen by Markov, gives
$\mathbb{E}N_{\rm pair}/(2M)\ge(1-2\delta)\log(1/(2\delta))-\log2$. Taking
$s=\sqrt{\mathbb{E}\tau_{\delta}}+1$ turns $\sqrt{\mathbb{E}\tau_{\delta}}\gtrsim\log(1/\delta)$
into the square.

\NOTEPAR{Scope: an ORACLE-ALLOCATION result. The tracked design is the surrogate's optimum
AT THE TRUE WORLD, held fixed. The true extrema choose which cells starve, so the true
parameter enters the allocation. The next paragraph removes this restriction. The stop is
the true-alternative GLR, so the cost measured is the design's, not a weaker stopping
rule's. The bound holds whether or not $\mathbb{E}\tau_{\delta}$ is finite.}

\leanbox{\lean{klq\_pair}, \lean{starved\_lower\_at}, \lean{starved\_cost\_log\_sq},
\lean{surrogate\_optimum\_cost}, \lean{surrogate\_optimum\_exists},
\lean{uniform\_surrogate\_cost}, \lean{surrogate\_vs\_ideal} in
\lean{LeanPaperA6SurrogateCostV1}; \lean{fCount\_le\_amgm}, \lean{fCount\_le\_sqrt} in
\lean{LeanPaperA6FixedDesignV1}; \lean{cnt\_eq\_at\_min}, \lean{ansEv\_alt\_leG},
\lean{ansEv\_true\_geG} in \lean{LeanPaperA7StoppedGenV1}}

\paragraph{The learner that re-solves the surrogate at its estimate.}
The deployable surrogate learner pays the same order. Each round it solves the frozen
surrogate at its current estimate $\hat\mu$ and tracks the result. At every estimate whose
extrema are the four anchors and which retains $A$, the surrogate game has value
$g(\hat\mu)^{2}/32$, with $g$ the frozen margin. There every exact optimum puts zero weight on
every middle cell, and the uniform-anchor design is one. Run the tracking learner with that
oracle, or with any oracle that returns an exact surrogate optimum at such estimates, with
the same forcing and the same stop. Then
$\mathbb{E}\tau_{\delta}\ge(KM/8)^{2}\log^{2}(1/\delta)$ for all $\delta<\delta_{0}$, and the ratio
to the ideal learner is unbounded. The proof adds one term to the one above: the rounds in
which some estimate is off by $1/16$ or more. Forcing gives every count order $\sqrt t$, so
after a deterministic time $s_{0}$ the learner's own confidence event places every estimate
within $1/16$ of the truth. That event fails at round $s$ with probability at most $s^{-4}$,
by Ville at the learner's statistic. So the expected number of misleading rounds is at most
$s_{0}+2$, at every horizon. It enters the pair's count additively and leaves the order
unchanged.

\NOTEPAR{Scope. The oracle sees only the estimate. The true world enters only as the law of
the observations and, in the analysis, as the centre of the $1/16$ ball. Where the estimate
misleads, the oracle is unconstrained (the implementation returns the uniform design), and
those rounds are paid for inside $s_{0}+2$. The stop is the true-alternative GLR, the family
is the separating one, and the forcing rate is $1/\sqrt{s+1}$. The threshold $\delta_{0}$ is
explicit but very small, because $s_{0}$ is. The statement is about the order in
$\log(1/\delta)$, not about any confidence level used in the experiments.
Figure~\ref{fig:adeploy} shows this learner, and the two others, under one two-answer
certificate at $\alpha=0.05$.}

\leanbox{\lean{est\_cost\_log\_sq}, \lean{surrogate\_learner\_cost},
\lean{exact\_oracle\_cost}, \lean{surrogate\_learner\_vs\_ideal}, \lean{est\_lower\_at} in
\lean{LeanPaperA6EstCostV1}; \lean{frozenGameAt\_eq}, \lean{starvation\_at},
\lean{surD\_exact}, \lean{optimal\_oracle\_middle\_zero}, \lean{close\_extreme} in
\lean{LeanPaperA6SurrogateAtV1}; \lean{radius\_eventually}, \lean{close\_of\_good},
\lean{badAll\_expect} in \lean{LeanPaperA6BadRoundsV1}; \lean{oCount\_le\_amgm},
\lean{oCount\_ge} in \lean{LeanPaperA6OracleLearnerV1}}

\section{Evidence}
\label{sec:evidence}

\begin{table}[t]
\caption{Evidence packages for this paper.}
\label{tab:evidence}
\centering
\footnotesize
\begin{tabular}{@{}llr@{}}
\toprule
Package & Script & Checks \\
\midrule
Frozen surrogate, both families & \texttt{EXP\_paperA\_frozen\_proxy} & see text \\
Certificate vs decision & \texttt{EXP\_unified\_oracle\_gap} & 7/7 \\
Support geometry, losses & \texttt{EXP\_unified\_geometry\_decision} & 31/31 \\
Optimised separation, both games & \texttt{EXP\_paperA\_separation\_games} & 11/11 \\
Repair criterion, rational witness & \texttt{EXP\_paperA\_repair\_barrier} & 15/15 \\
Learner: stopping with a certificate & \texttt{EXP\_paperA7\_learner} & 14/14 \\
Track-and-stop, window accumulation & \texttt{EXP\_paperA7\_track\_and\_stop} & 6/6 \\
Does the rule learn the answer? & \texttt{EXP\_paperA\_swap} & 5/5 \\
Grid rate and the stopping chain & \texttt{EXP\_paperA7\_stopping\_chain} & 10/10 \\
Allocation mechanism, one certificate & \texttt{EXP\_paperA\_mechanism} & 5/5 \\
Mechanism data against the theorems & \texttt{FIG\_paperA\_mechanism} & 7/7 \\
\bottomrule
\end{tabular}
\end{table}

\begin{figure}[t]
\centering
\includegraphics[width=\columnwidth]{paperA7_stopping_chain.pdf}
\caption{Left: at $\Delta=0.30$ the rate $u\Delta-u^{2}/2$ is negative at every non-zero
point of a fixed mesh $d=1$, while a mesh $m=0.05$ lands above $\Delta^{2}/2-m^{2}/2$ ---
the loss is the squared miss, so the fixed mesh is the whole obstruction. Right: at the
window $n=\lceil\sqrt t\,\rceil$ the failure mass $1/(n-1)^{3}$ sits under the proved
$8t^{-3/2}$, which is summable; the per-step confidence input is not simulated.}
\label{fig:a7chain}
\end{figure}
\begin{figure}[t]
\centering
\includegraphics[width=\columnwidth]{paperA7_learner.pdf}
\caption{Left: stopping cost per unit of confidence for the three arms, against the
asymptotic \emph{bracket} for $1/\Gamma_{\rm comp}$. Right: every stopped run returned the
correct answer. Finite-$\alpha$ measurements; no ordering is asserted.}
\label{fig:learner}
\end{figure}

\begin{figure}[t]
\centering
\includegraphics[width=\columnwidth]{paperA_separation_games.pdf}
\caption{\textbf{A different construction from Theorem~\ref{thm:sepcentred}, and not a
picture of it.} The two optimised ratios against $M$ on the $r$-pair/frozen-proxy family,
each with the bound Theorem~\ref{thm:sep} proves; that family's surrogate constant is
$\Gamma_{\rm proxy}=1/512$, whereas the centred construction's is
$\Gamma_{\rm frozen}=9/512$. The two must not be relabelled into each other. Cutting planes
bracket both games; the brackets, not a single number, are what is plotted.}
\label{fig:sepgames}
\end{figure}

\begin{figure}[t]
\centering
\includegraphics[width=\columnwidth]{paperA_repair_barrier.pdf}
\caption{Where the repair criterion fires, and the shaded band where the proved ceiling
$64/M$ exceeds $1$ and therefore says nothing.}
\label{fig:repair}
\end{figure}

\begin{figure}[t]
\centering
\includegraphics[width=\columnwidth]{paperA_frozen_proxy.pdf}
\caption{The surrogate misprices in both directions, against the predicted rates.}
\label{fig:proxy}
\end{figure}

\begin{figure}[t]
\centering
\includegraphics[width=0.78\columnwidth]{paperA_oracle_gap.pdf}
\caption{Certifying the margin against making the decision.}
\label{fig:margin}
\end{figure}

\begin{figure}[t]
\centering
\includegraphics[width=\columnwidth]{paperA_certified.pdf}
\caption{A-F1: the centred separation from a numerical true-game computation that includes
centring and every alternative region, certified in exact rational arithmetic
(\texttt{EXP\_paperA\_certified\_separation}, 6/6; figure 3/3). The upper end comes from the dual of
the outer LP over exactly-verified centred alternatives. The lower end is the minimum over
all $|B|(|B|-1)$ regions of a Lagrangian dual value at a design with every weight positive.
Relative widths are at most $4\times10^{-4}$: $\Gamma_{\rm true}/\Gamma_{\rm proxy}\in[0.8372,0.8374]$,
$[0.7214,0.7216]$, $[0.5664,0.5667]$, $[0.3982,0.3982]$ at $r=1,2,4,8$, inside the proved
bracket of Theorem~\ref{thm:sepcentred} (band). $\Gamma_{\rm proxy}=9/512$ is the theorem, not a
computation.}
\label{fig:acert}
\end{figure}

\begin{figure*}[t]
\centering
\includegraphics[width=0.92\textwidth]{paperA_mechanism.pdf}
\caption{The centred separation and the mechanism, from one stored panel ($M=1$, $\alpha=0.05$,
24 runs per rule and $r$, budget $B=6\times10^{4}$); nothing was re-solved or re-sampled.
Left: the band is the proved bracket $4/(4+2rM)\le\Gamma_{\rm true}/\Gamma_{\rm proxy}\le
64/(9rM)$. Each interval runs from the best attained LOWER certificate on the true value
(floating point; the strict centred alternatives are relaxed to a union of polyhedra, which
can only lower the infimum) up to the theorem's upper bound. No numerical upper bound is
claimed here; Figure~\ref{fig:acert} gives a two-sided certified interval with centring. At $r=8$ the whole interval lies below $1$. Centre: the share of queries on the
middle cells. The frozen proxy's optimum puts nothing there, and its runs sample the middles
only through the forced exploration. Right: $\mathbb{E}[\min(\tau,B)]$ with its standard
error. All proxy runs are censored at $B$, which does not show that they never stop. SCOPE:
the stopping test certifies answer $A$ only, and the geometry rule optimises that same
certificate. The answer at the truth is therefore fixed in advance: this is an ALLOCATION
diagnostic under one certificate, not a $\delta$-correct learner (for that question see
Figure~\ref{fig:aswap}). Every stored number drawn was checked against the proved formula
beside it before drawing (\texttt{FIG\_paperA\_mechanism}, 7/7).}
\label{fig:amech}
\end{figure*}

\begin{figure*}[t]
\centering
\includegraphics[width=0.92\textwidth]{paperA_budget.pdf}
\caption{The same stored panel on the axis A-F2 asks for: budget, not a single $B$. Nothing
was re-solved or re-sampled. Left: the certificate rate $\Pr(\tau\le b)$ from the per-run
stopping times. The surrogate-optimum rule certifies in none of its $24$ runs at any
$b\le B$, for every $r$. Centre: the restricted mean $\mathbb{E}[\min(\tau,b)]$, which is
exact for every $b\le B$ because a censored run has $\min(\tau,b)=b$. The surrogate-optimum
rule lies on the diagonal. Right: $N_{\rm mid}(t)/t$ along the FIRST run of each rule only,
for $r=1$ and $r=8$. The black curves are the pathwise forcing bound
$2r\,(2+\sum_{n\le s<t}s^{-1/2}/n)$ of \lean{fCount\_le\_quota}, transcribed for the code's
forcing $1/\sqrt t$. The surrogate-optimum run stays below it, within one query at $r=1$:
forcing is its only source of middle queries. Line style marks $r$ in every panel. SCOPE: the
surrogate-optimum rule is the Lean \lean{fixLearner} at \lean{wUnifC}, up to a one-step
offset in the forcing. Its stop is the answer-$A$ certificate of Figure~\ref{fig:amech}, not
the GLR stop of the theorem, and $\alpha=0.05$ lies outside the theorem's range
$\delta<\delta_{0}$ (about $1.7\times10^{-4}$ at $r=1$, $M=1$). So the panel is consistent
with the $\log^{2}$ separation in direction and is not a test of it
(\texttt{FIG\_paperA\_budget}, 5/5 checks).}
\label{fig:abudget}
\end{figure*}

\begin{figure*}[t]
\centering
\includegraphics[width=0.92\textwidth]{paperA_deployable.pdf}
\caption{Deployable policies under ONE two-answer certificate: the true answer no longer
enters the stop. Every rule is a tracking learner with forcing $\min(1,1/\sqrt t)$ whose
oracle sees only the estimate. \emph{Surrogate at estimate} re-solves the frozen surrogate
at the current estimate every round ($w_{e}\propto\sigma_{e}$ on the estimate's four frozen
extremum cells, checked to be its exact optimum). \emph{True game at estimate} ascends the
certificate of the answer the estimate retains. Each run stops when the certificate of that
answer clears the threshold, and returns it. Left, centre: the certificate rate by budget,
on the base world (answer $A$) and on the swapped world (answer $B$), for $r=1,2,4,8$.
Right: $\mathbb{E}[\min(\tau,B)]$, $B=6\times10^{4}$, $\alpha=0.05$, 24 runs per cell. Of 505
stopped runs, none returned the wrong answer. On the base world the surrogate learner
certifies in 22 of 24 runs at $r=1$, 3 at $r=2$ and none at $r=4,8$, while the other two
rules certify every run. On the swapped world, where the surrogate's frozen cells are the
informative ones, it certifies every run, and is the fastest rule except at $r=2$, where it
and the true-game rule are within $1\%$. Along every run the middle counts obey the forcing
bound of \lean{oCount\_le\_bad}, with minimum slack $0.33$. SCOPE: the certificate is a valid
lower bound on the likelihood ratio, not the exact infimum of the theorem's stop, and
$\alpha=0.05$ is far above the theorem's $\delta_{0}$, so the panel shows these rules'
behaviour. It is not a test of the $\log^{2}$ order (\texttt{EXP\_paperA\_deployable}, 6/6;
\texttt{FIG\_paperA\_deployable}, 4/4).}
\label{fig:adeploy}
\end{figure*}

\section{Limitations}
\label{sec:limits}

\begin{enumerate}
\item[1)] \textbf{The two families are not yet realised in the model.} They are stated for a
range oracle with per-cell observations; the bridge of Section~\ref{sec:proxy} is owed.
\item[2)] \textbf{The certificates are numerical, not exact.} Every one of the $81$ regions
of every cell passes the KKT check, degenerate ones included: the multipliers are recovered
by \emph{non-negative} least squares, since on a region with linearly dependent active
constraints---one pinning the same index as both maximum and minimum---the minimum-norm
multipliers can be negative while admissible ones exist, and it is existence that
Theorem~\ref{thm:kkt} asks for. The residuals are floating-point, so this certifies
the decomposition to solver precision and not in exact arithmetic. A region that failed
would be reported and would fail the run's self-check, never dropped.
\item[3)] \textbf{The implemented learner's rate is conditional.} For the ideal learner of
Section~\ref{sec:adaptive} the stopping chain is complete. The per-step confidence guarantee
$\Pr(\text{fails at }s)\le s^{-4}$ is Ville's inequality at the learner's own statistic, and
the rate, finiteness and correctness are proved. For the implementation of
Section~\ref{sec:learner} (finite precision, a bracketing oracle, periodic refresh) the same
deterministic chain still takes the hypotheses named there. That learner is reported only as
a finite-panel measurement. The assembly step
$\mathbb{E}\tau\le t_{0}+\sum_{t\ge t_{0}}\Pr(\tau>t)$ is proved with no term discarded
(\lean{expectation\_le\_split} in \lean{LeanPaperClosuresV1}).
\item[4)] \textbf{The model is fixed, not derived} from the corpus definitions.
\item[5)] \textbf{Finitely many support points and competitors}, and nothing bounds the
number of regions, which is what decides whether enumeration is practical.
\item[6)] \textbf{The sample-cost separation is for one family and one forcing rate.} It
holds for the learner that re-solves the surrogate at its estimate, with forcing
$1/\sqrt{s+1}$ and the true-alternative GLR stop, on the separating family, for
$\delta<\delta_{0}$ with an explicit but very small $\delta_{0}$. Other forcing rates, other
stops and other families are not covered.
\end{enumerate}

\section{Conclusion}

The alternative set of a support-aware compression decision can be pinned exactly rather
than approximated, its regional minima and its design can both be certified rather than
searched, and the surrogate that freezes the extrema is wrong in both directions by
unbounded factors. Separately, certifying the margin that would justify a decision is up to
three orders of magnitude more expensive than the decision, and past a threshold impossible
while the decision stays easy. For the ideal learner the stopping chain is complete: it attains $1/\Gamma_{\rm true}$, and
the learner that re-solves the surrogate at its estimate instead costs
$\Omega(\log^{2}(1/\delta))$ samples. What remains owed is the bridge realising the two
counterexample families inside the linear model, and the rate of the implemented
finite-precision learner.

\section*{Companion material}
Two files accompany this paper and are part of the submission record.
\texttt{PROOFS.md} gives the central arguments as readable mathematics with no Lean
required --- setup, claim, proof, and in the same breath what the claim does not
cover. \texttt{DOSSIER.md} carries the full state: every audited theorem, every
measured number with its artifact and provenance status, the open list, and the
claims retracted during preparation, recorded rather than silently edited.
\texttt{CLAIMS.md} is the one-row-per-headline table of claim, final theorem,
remaining hypotheses and evidence.

\begin{thebibliography}{00}

\bibitem{schutte2026} N. Schutte et al., ``Sufficient decision proxies for
decision-focused learning,'' arXiv:2505.03953, 2026.

\bibitem{tosh2022} C. Tosh et al., ``Targeted active learning for probabilistic models,''
arXiv:2210.12122, 2022.


\bibitem{overstall2022} A. M. Overstall, ``Properties of Fisher information gain for
Bayesian design of experiments,'' Univ. of Southampton, 2022. [Online]. Available:
https://eprints.soton.ac.uk/452554/

\bibitem{filstroff2021} L. Filstroff, I. Sundin, P. Mikkola, A. Tiulpin, J. Kylm\"{a}oja,
and S. Kaski, ``Targeted active learning for Bayesian decision-making,'' arXiv:2106.04193,
2021.

\bibitem{tang2026} Y. Tang, S. Sloman, and S. Kaski, ``Error amplification under model
misspecification,'' in \emph{Proc. Machine Learning Research}, vol.~300, 2026.

\bibitem{bendavid2026} S. Ben-David and R. Kothari, ``A total function with randomized
query complexity quadratically below its certificate complexity,'' arXiv:2609.15063, 2026.

\bibitem{wan2026} Z. Wan, Y. Liu, P. Grigas, and Z.-J. M. Shen, ``Decision-focused
sequential experimental design,'' arXiv:2602.05340, 2026.


\bibitem{degenne2019} R. Degenne and W. M. Koolen, ``Pure exploration with multiple correct
answers,'' in \emph{Advances in Neural Information Processing Systems 32}, 2019,
pp.~14591--14600.

\bibitem{wangyou2026} Z. Wang and K. You, ``Information games with multiple correct answers
and piecewise-convex alternative sets,'' arXiv:2606.21889, 2026.

\bibitem{bennouna2026} A. Bennouna et al., ``Data informativeness in linear optimization,''
arXiv:2602.15365, 2026.

\bibitem{boyd2004} S. Boyd and L. Vandenberghe, \emph{Convex Optimization}. Cambridge,
U.K.: Cambridge Univ. Press, 2004.

\bibitem{pukelsheim2006} F. Pukelsheim, \emph{Optimal Design of Experiments}. Philadelphia,
PA: SIAM Classics in Applied Mathematics, 2006.

\bibitem{nitinawarat2013} S. Nitinawarat and V. V. Veeravalli, ``Controlled sensing for
sequential multihypothesis testing with controlled Markovian observations and non-uniform
control cost,'' \emph{Sequential Analysis}, vol.~34, no.~1, pp.~1--24, 2015.

\bibitem{mathlib2020} The mathlib Community, ``The Lean mathematical library,'' in
\emph{Proc. 9th ACM SIGPLAN Int. Conf. Certified Programs and Proofs (CPP)}, 2020,
pp.~367--381.

\bibitem{demoura2021} L. de Moura and S. Ullrich, ``The Lean 4 theorem prover and
programming language,'' in \emph{Proc. 28th Int. Conf. Automated Deduction (CADE)}, 2021,
pp.~625--635.

\end{thebibliography}

\end{document}
